% This work is licensed under the Creative Commons
% Attribution-NonCommercial-ShareAlike 4.0 International License. To view a copy
% of this license, visit http://creativecommons.org/licenses/by-nc-sa/4.0/ or
% send a letter to Creative Commons, PO Box 1866, Mountain View, CA 94042, USA.
% (c) Eric Kunze, 2020
%%%%%%%%%%%%%%%%%%%%%%%%%%%%%%%%%%%%%%%%%%%%%%%%%%%%%%%%%%%%%%%%%%%%%%%%%%%%%%%
% Template for lecture notes at TU Dresden.
%%%%%%%%%%%%%%%%%%%%%%%%%%%%%%%%%%%%%%%%%%%%%%%%%%%%%%%%%%%%%%%%%%%%%%%%%%%%%%%

\documentclass[ngerman, a4paper, 11pt]{report}

\usepackage[ngerman]{babel}
\usepackage[top=2.5cm,bottom=2.5cm,left=2.5cm,right=2.5cm]{geometry}

\usepackage{parskip}    % split paragraphs by vspace instead of intendations
\usepackage[onehalfspacing]{setspace} % increase row-space

\usepackage{xifthen}

%%%%%%%%%%%%%%%%%%%%%%%%%%%%%%%%%%%%%%%%%%%%%%%%%%%%%%%%%%%%%%%%%%%%%%%%%%%%%%%
% SCHRIFTEN
\usepackage[utf8]{inputenc}
\usepackage{chngcntr}
\usepackage{eufrak}

\usepackage{lmodern}
\usepackage[normalem]{ulem}

\usepackage[autostyle=true,english=british]{csquotes}

% new font OpenSans
\usepackage[scale=1]{opensans}
\newcommand*{\osfamily}{\fontfamily{fos}\selectfont}
\DeclareTextFontCommand{\textos}{\osfamily}

%%%%%%%%%%%%%%%%%%%%%%%%%%%%%%%%%%%%%%%%%%%%%%%%%%%%%%%%%%%%%%%%%%%%%%%%%%%%%%%
% TABLE OF CONTENTS
\usepackage{tocloft}    % customize toc
\renewcommand{\cfttoctitlefont}{\titlefont\Huge\bfseries}
\setcounter{tocdepth}{1}
\renewcommand{\cftbeforetoctitleskip}{0pt} % change appearence of heading of toc: 0 space above, bold, titlefont, huge toc-heading
\renewcommand{\cftchapnumwidth}{2em} % change indentations due to width of capital roman numbers
\renewcommand{\cftsecindent}{2em}
\renewcommand{\cftsecnumwidth}{2em}
\renewcommand{\cftsubsecindent}{4em}

%%%%%%%%%%%%%%%%%%%%%%%%%%%%%%%%%%%%%%%%%%%%%%%%%%%%%%%%%%%%%%%%%%%%%%%%%%%%%%%
% GRAPHICS
\usepackage{graphicx}
\usepackage[table,dvipsnames]{tudscrcolor} % package xcolor loaded automatically
\usepackage[font=small,labelfont=bf]{caption} % captions of non-floated figures

\usepackage{tcolorbox}
\usepackage{tikz}

\usepackage[table,dvipsnames]{tudscrcolor} % package xcolor loaded automatically

% tabulars
\usepackage{tabularx}   % tabularx-environment (explicitly set width of columns)
\usepackage{multirow}
\usepackage{booktabs}   % improved rules

%%%%%%%%%%%%%%%%%%%%%%%%%%%%%%%%%%%%%%%%%%%%%%%%%%%%%%%%%%%%%%%%%%%%%%%%%%%%%%%
% TITLE STYLES
\usepackage{titlesec}   % change title headings look
\usepackage{relsize}    % relative font size (smaller[i], larger[i], ...)

\newcommand{\titlefont}{\osfamily}
\newcommand{\chaptersize}{\huge}
\newcommand{\sectionsize}{\LARGE}

\renewcommand{\thepart}{\Alph{part}}

% \titleformat{<command>}[<shape>]{<format>}{<label>}{<sep>}{<before-code>}[<after-code>]
% \titlespacing*{<command>}{<left>}{<before-sep>}{<after-sep>}[<right-sep>]

%%%%%%%%% framed chapter
\titleformat{\chapter}[frame]{\bfseries\titlefont\color{cddarkblue}}{\enskip \chaptersize \smaller \chaptername\;\thechapter \enskip}{10pt}{\chaptersize\centering\MakeUppercase}%
\titlespacing{\chapter}{0pt}{0pt}{10pt}%

%%%%%%%%% chapter.section
\counterwithin{section}{chapter}%
\titleformat*{\section}{\bfseries\titlefont\sectionsize}%{\thesection}{8pt}{}%
\titlespacing{\section}{0pt}{10pt}{5pt}
\titleformat*{\subsection}{\bfseries\titlefont\sectionsize\smaller}

\pretocmd{\chapter}{\setcounter{section}{0}}{}{} % automatic reset of section after chapter ended
\pretocmd{\section}{\setcounter{equation}{0}}{}{} % automatic reset of equation counter in each section

%%%%%%%%%%%%%%%%%%%%%%%%%%%%%%%%%%%%%%%%%%%%%%%%%%%%%%%%%%%%%%%%%%%%%%%%%%%%%%%
% ENUMERATIONS
\usepackage{enumerate}
\usepackage[inline]{enumitem}       % customize label

\renewcommand{\labelitemi}{\raisebox{1pt}{\scalebox{.4}{$\blacksquare$}}}
\renewcommand{\labelitemii}{$\vartriangleright$}
\renewcommand{\labelitemiii}{--}
% Variantionen des Dreiecks als Aufzählungszeichen $\blacktriangleright$ / $\vartriangleright$ / $\triangleright$

\renewcommand{\labelenumi}{(\arabic{enumi})}
\renewcommand{\labelenumii}{\alph{enumii}.}
\renewcommand{\labelenumiii}{\roman{enumiii}.}

%%%%%%%%%%%%%%%%%%%%%%%%%%%%%%%%%%%%%%%%%%%%%%%%%%%%%%%%%%%%%%%%%%%%%%%%%%%%%%%
% PAGE STYLES
\newcommand*{\rightinfo}{Vorlesung \enquote{Mathematische Statistik} bei Prof. Dr. Ferger im Wintersemester 2020/21}

\usepackage{fancyhdr}   % customize header / footer

\fancypagestyle{myStyle}{%
    \fancyhf{} %
    \fancyfoot[C]{\thepage} %
    \renewcommand{\headrulewidth}{0pt}     % Line at the header invisible
    \renewcommand{\footrulewidth}{0pt}     % Line at the footer visible
    \fancyhead[C]{\textcolor{gray}\leftmark} %
    \fancyhead[R]{%
        \begin{tikzpicture}[overlay,remember picture]
            \node [
                fill=none,  % Farbe des Randstreifens
                text=gray,  % Textfarbe
                font=\osfamily\normalsize,  % Einstellungen für die Schrift
                inner xsep=\footskip,       % Abstand des Textes von unten
                % maximale Textbreite = Papierhöhe - 2*Abstand des Textes von unten:
                text width={\dimexpr\paperheight-2\footskip\relax},
                align=center,
                minimum height=7mm,% Breite des Randstreifens
                anchor=south west,
                rotate=90
            ] at ($(current page.south east)+(-10mm,0mm)$)
            {\rightinfo};
        \end{tikzpicture}%
    }
}
\fancypagestyle{rightinfo}{%
    \fancyhf{} %
    \fancyfoot[C]{\thepage} %
    \renewcommand{\headrulewidth}{0pt}     % Line at the header invisible
    \renewcommand{\footrulewidth}{0pt}     % Line at the footer visible
    \fancyhead[R]{%
        \begin{tikzpicture}[overlay,remember picture]
            \node [
                fill=none,  % Farbe des Randstreifens
                text=gray,  % Textfarbe
                font=\sffamily\normalsize,  % Einstellungen für die Schrift
                inner xsep=\footskip,       % Abstand des Textes von unten
                % maximale Textbreite = Papierhöhe - 2*Abstand des Textes von unten:
                text width={\dimexpr\paperheight-2\footskip\relax},
                align=center,
                minimum height=7mm,% Breite des Randstreifens
                anchor=south west,
                rotate=90
            ] at ($(current page.south east)+(-10mm,0mm)$)
            {\rightinfo};
        \end{tikzpicture}%
    }
}
% changes pagestyle on first page of each chapter; instead of empty page the normal footer is printed
\patchcmd{\chapter}{\thispagestyle{plain}}{\thispagestyle{rightinfo}}{}{}

\pagestyle{myStyle}
\pagenumbering{arabic}

%%%%%%%%%%%%%%%%%%%%%%%%%%%%%%%%%%%%%%%%%%%%%%%%%%%%%%%%%%%%%%%%%%%%%%%%%%%%%%%
% TITLEPAGE
\newcommand{\makeTUtitle}[1][]{%
    \begin{titlepage}
        \pagecolor{cddarkblue!90}
        \color{white}

        \raggedright
        \fosfamily
        \setlength{\parindent}{0pt}

        \hspace{-18.6mm} %
        \includegraphics[scale=0.6]{TUD-white.pdf} \\
        %       \includegraphics[scale=0.6]{TUD-blue.pdf} \\

        \vspace{3mm}
        \begin{tabular}{m{\textwidth}}
            \hline
            \hspace{-4pt}\small{\textbf{Fakultät Mathematik} Institut für Stochastik, Professur für Statistik} \\
            \hline
        \end{tabular} \\

        \vspace{5cm}
        {\Huge\bfseries \MakeUppercase Mathematische Statistik \par}
        %
        \ifthenelse{\isempty{#1}}{}{
            \vspace{0.5cm}%
            {\Large \itshape #1} \\%
        }
        %
        \vspace{1.5cm}
        \textbf{{\Large Prof. Dr. Dietmar Ferger}} \par
        \vspace{0.5cm}
        {\large Wintersemester 2020/21}

        \vfill
        \begin{tabular}{lll}
            Mitschrift  & : & Eric Kunze \\
            E-Mail      & : & \url{eric.kunze@mailbox.tu-dresden.de} \\
        \end{tabular}
    \end{titlepage}
    \nopagecolor
}

%%%%%%%%%%%%%%%%%%%%%%%%%%%%%%%%%%%%%%%%%%%%%%%%%%%%%%%%%%%%%%%%%%%%%%%%%%%%%%%
% MATH ENVIRONMENTS

\usepackage{../mathoperatorsMathTUD}

\usepackage[ntheorem,framemethod=TikZ]{mdframed}
\usepackage{amsmath,amssymb,amsfonts,mathtools}
\usepackage[amsmath,thmmarks,framed]{ntheorem}

\newcommand{\skiparound}{10pt}

\newcounter{thmcount}
\newcounter{defcount}
\counterwithin{thmcount}{chapter}
\counterwithin{defcount}{chapter}

% possible settings:
% style=boxedtheorem,%
% % backgroundcolor=cdpurple!5,%
% % linecolor=cdpurple,%
% innertopmargin=\boxskip,%
% innerbottommargin=\boxskip,%
% % innerleftmargin=4pt,%
% % outerlinewidth=5pt,%
% % topline=false,%
% % rightline=true,%
% % leftline=false,%
% % bottomline=false,%
% innertopmargin=3pt,%
% innerbottommargin=3pt,%
% % leftmargin=-4pt,%
% rightmargin=-10pt,%
% frametitlefont=\osfamily\bfseries,%
% skipabove=5pt,%
% skipbelow=5pt,%
% % skipabove=6pt,%
% % skipbelow=6pt,%
% % nobreak,%

\theorempreskip{\skiparound}
\theorempostskip{\skiparound}

\theoremstyle{plain}
\theoremseparator{}
\theoremheaderfont{\osfamily \bfseries \upshape}
\theorembodyfont{}

% DEFINTION
\newmdtheoremenv[%
	backgroundcolor=cdorange!10,%
	linecolor=cdorange,%
	innerleftmargin=4pt,%
	outerlinewidth=5pt,%
	topline=false,%
	rightline=true,%
	leftline=false,%
	bottomline=false,%
	leftmargin=-4pt,%
	skipabove=6pt,%
	skipbelow=6pt,%
	nobreak,%
]{definition}[thmcount]{Definition}

\theoremstyle{plain}
\theoremseparator{}

% THEOREM
\newmdtheoremenv[%
	backgroundcolor=cdblue!10,%
	linecolor=cddarkblue!90,%
	innerleftmargin=4pt,%
	outerlinewidth=5pt,%
	topline=false,%
	rightline=true,%
	leftline=false,%
	bottomline=false,%
	leftmargin=-4pt,%
	skipabove=6pt,%
	skipbelow=6pt,%
	nobreak,%
]{theorem}[thmcount]{Theorem}

% SATZ
\newmdtheoremenv[%
	backgroundcolor=cdblue!10,%
	linecolor=cddarkblue!90,%
	innerleftmargin=4pt,%
	outerlinewidth=5pt,%
	topline=false,%
	rightline=true,%
	leftline=false,%
	bottomline=false,%
	leftmargin=-4pt,%
	skipabove=6pt,%
	skipbelow=6pt,%
]{satz}[thmcount]{Satz}

% LEMMA
\newmdtheoremenv[%
	backgroundcolor=cdindigo!4,%
	linecolor=cdindigo!50,%
	innerleftmargin=4pt,%
	outerlinewidth=5pt,%
	topline=false,%
	rightline=true,%
	leftline=false,%
	bottomline=false,%
	leftmargin=-4pt,%
	skipabove=6pt,%
	skipbelow=6pt,%
	nobreak,%
]{lemma}[thmcount]{Lemma}

% PROPOSITION
\newmdtheoremenv[%
	backgroundcolor=cdindigo!4,%
	linecolor=cdindigo!50,%
	innerleftmargin=4pt,%
	outerlinewidth=5pt,%
	topline=false,%
	rightline=true,%
	leftline=false,%
	bottomline=false,%
	leftmargin=-4pt,%
	skipabove=6pt,%
	skipbelow=6pt,%
	nobreak,%
]{proposition}[thmcount]{Proposition}

% KOROLLAR
\newmdtheoremenv[%
	backgroundcolor=cdpurple!5,%
	linecolor=cdpurple,%
	innerleftmargin=4pt,%
	outerlinewidth=5pt,%
	topline=false,%
	rightline=true,%
	leftline=false,%
	bottomline=false,%
	leftmargin=-4pt,%
	skipabove=6pt,%
	skipbelow=6pt,%
	nobreak,%
]{korollar}[thmcount]{Korollar}

\theoremstyle{plain}
\theoremheaderfont{\osfamily \bfseries}
\theorembodyfont{}

\newtheorem{beispiel}[thmcount]{Beispiel}
\newtheorem{erinnerung}[thmcount]{Erinnerung}
\newtheorem{wiederholung}[thmcount]{Wiederholung}
\newtheorem{bemerkung}[thmcount]{Bemerkung}

% theorems without numbers
\theoremstyle{nonumberplain}
\theoremheaderfont{\osfamily \bfseries}
\theorembodyfont{}
\theoremseparator{.}

\newmdtheoremenv[%
	backgroundcolor=cdorange!10,%
	linecolor=cdorange,%
	innerleftmargin=4pt,%
	outerlinewidth=5pt,%
	topline=false,%
	rightline=true,%
	leftline=false,%
	bottomline=false,%
	leftmargin=-4pt,%
	skipabove=6pt,%
	skipbelow=6pt,%
	nobreak,%
]{*definition}{Definition}

\newtheorem{*interpretation}{Interpretation}
\newtheorem{*bemerkung}{Bemerkung}
\newtheorem{*beispiel}{Beispiel}
\newtheorem{*erinnerung}{Erinnerung}
\newtheorem{*notation}{Notation}

% >> proof
\makeatletter
\newtheoremstyle{proofstyle}%
{\item[\hskip\labelsep {\theorem@headerfont ##1}\theorem@separator]}%
{\item[\hskip\labelsep {\theorem@headerfont ##1}\ (##3)\theorem@separator]}
\makeatother

\theoremstyle{proofstyle}
\theoremheaderfont{\normalsize\slshape}
\theorembodyfont{}
\theoremseparator{.}
\theorempreskip{5pt}
\theorempostskip{5pt}
\theoremsymbol{$\square$}
\newtheorem{proof}{Beweis}

% >> Equivalences
\newcommand{\hinrichtung}{\item[\bfseries ($\boldsymbol{\Rightarrow}$)]}
\newcommand{\rueckrichtung}{\item[\bfseries ($\boldsymbol{\Leftarrow}$)]}
\newenvironment{equivalence}[1][]{%
        \ifthenelse{\isempty{#1}}%
        {\begin{description}} % no optional argument
            {\begin{description}[topsep=-\parskip]} % any optional argument
    }{ \end{description} }%

    % >> inductions (needs description environment)
    \newcommand{\ianfang}[1][]{%
        \ifthenelse{\isempty{#1}}{%
            \item[\textbf{(IA)}] % no parameter
        }{  \item[\textbf{(IA)}] {#1 :}  } % parameter exists
    }

    \newcommand{\ivorraussetzung}{\item[\bfseries (IV)]}

    \newcommand{\ischritt}[1][]{%
        \ifthenelse{\isempty{#1}}{%

            \item[\textbf{(IS)}] % no parameter
        }{  \item[\textbf{(IS)}] {#1 :} } % parameter exists
    }

%%%%%%%%%%%%%%%%%%%%%%%%%%%%%%%%%%%%%%%%%%%%%%%%%%%%%%%%%%%%%%%%%%%%%%%%%%%%%%%
% HYPERLINKS
\usepackage[unicode,bookmarks=true]{hyperref}
\hypersetup{
    % pdfborder={0 0 0}         % no boxed around links
    pdfborderstyle={/S/U/W 1},  % underlining insteas of boxes
    linkbordercolor=cdblue,
    urlbordercolor=cdblue
    %   colorlinks,
    %   citecolor=black,
    %   filecolor=cddarkblue!80,
    %   linkcolor=black,
    %   urlcolor=cddarkblue!80
}
\usepackage{bookmark}       % pdf-bookmarks

\usepackage{cleveref}
\crefname{theorem}{Theorem}{Theoreme}
\crefname{satz}{Satz}{Sätze}
\crefname{lemma}{Lemma}{Lemmata}
\crefname{proposition}{Proposition}{Propositionen}
\crefname{folgerung}{Folgerung}{Folgerungen}
\crefname{korollar}{Korollar}{Korollare}
\crefname{definition}{Definition}{Definitionen}
\crefname{bemerkung}{Bemerkung}{Bemerkungen}
\crefname{beispiel}{Beispiel}{Beispiele}
\crefname{erinnerung}{Erinnerung}{Erinnerungen}

%%%%%%%%%%%%%%%%%%%%%%%%%%%%%%%%%%%%%%%%%%%%%%%%%%%%%%%%%%%%%%%%%%%%%%%%%%%%%%%
% HIGHLIGHTING
\newcommand{\begriff}[1]{\textbf{#1}}
\newcommand{\person}[1]{\textsc{#1}}

%%%%%%%%%%%%%%%%%%%%%%%%%%%%%%%%%%%%%%%%%%%%%%%%%%%%%%%%%%%%%%%%%%%%%%%%%%%%%%%
% ADDITIONAL COMMANDS

\usepackage[numbers, square]{natbib}

\renewcommand{\dx}{\mathrm{d}x}
\newcommand{\A}{\mathcal{A}}
\newcommand{\B}{\mathcal{B}}
\renewcommand{\S}{\mathcal{S}}
\renewcommand{\G}{\mathcal{G}}
\renewcommand{\L}{\mathcal{L}}
\renewcommand{\F}{\mathcal{F}}
\newcommand{\ntoinf}{\overset{n \to \infty}{\longrightarrow}}
\newcommand{\satzende}{\enskip \text{.}}
\newcommand{\komma}{\enskip \text{,}}

%%%%%%%%%%%%%%%%%%%%%%%%%%%%%%%%%%%%%%%%%%%%%%%%%%%%%%%%%%%%%%%%%%%%%%%%%%%%%%%

\begin{document}
    \makeTUtitle
    
    ~ \vfill
    \begin{center}
    	\fosfamily \itshape
    	Dieses Skript basiert auf den Mitschriften von Willi Sontopski und Felix Hilsky aus den vorherigen Semestern sowie den Bemerkungen und Korrekturen von Prof. Ferger im Wintersemester 2020/21. 
    \end{center}
	\vfill
	
	\pagebreak

    \tableofcontents

    \chapter{Median}
    \label{chapter_1_median}
    Viele Schätzer in der Statistik sind definiert als Minimal- oder Maximalstelle von bestimmten \begriff{Kriteriumsfunktionen}, z.\,B. der \begriff{Maximum-Likelihood-Schätzer (MLS)} oder \begriff{Mi\-ni\-mum-Quadrat-Schätzer (MQS, KQS)} oder \begriff{Bayes-Schätzer}. Allgemein nennt man solche Schätzer \begriff{M-Schätzer}.

\begin{beispiel}[Maximum-Likelyhood-Schätzer]
	Gegeben seien $X_1, ..., X_n$ iid. $\sim f_\theta$. Dann ist $\hat{\theta}_n$ die Maximumstelle von der Funktion $l \colon \theta \mapsto \sum_{i = 0}^n \log f_{\theta}(X_i)$ (sogenannte Log-Likelyhood-Funktion). Dabei kommen alle $\theta$ aus einer möglichen Menge $\Theta$ in Frage.
\end{beispiel}

Ziel: Untersuchung des asymptotischen Verhaltens ($n \to \infty$) von M-Schätzern über einen funktionalen Ansatz. Als Beispiel betrachten wir nun den Median

    % This work is licensed under the Creative Commons
% Attribution-NonCommercial-ShareAlike 4.0 International License. To view a copy
% of this license, visit http://creativecommons.org/licenses/by-nc-sa/4.0/ or
% send a letter to Creative Commons, PO Box 1866, Mountain View, CA 94042, USA.

Sei $X \colon (\Omega,\A,\P) \to \R$ eine reelle Zufallsvariable mit Verteilungsfunktion $F_X\colon \R \to [0,1]$, d.\,h.
\begin{equation*}
	F_X(x) \defeq \P(X \le x) 
	\defeq \P(\menge{\omega \in \Omega : X(\omega) \le x})
	= \P\brackets{X^{-1}\brackets{(-\infty, x]}},
\end{equation*}
also $X \sim F_X$. Definiere
\begin{align}
	Y(t) &\defeq \E(\abs{X-t}) \label{eq: DefY} \\ \nonumber
	&= \int_\Omega \abs{X(\omega)-t} \, \P(\diff\omega) \\ \nonumber
	\overset{\cref{eqTrafoR}}&{=} \int_\R \abs{x-t} \, (\P \circ X^{-1})(\diff x) \\ \nonumber
	&= \int_\R \abs{x-t} \, F_X(\diff x) \qquad \forall t \in \R \\
	\intertext{und}
	m &\in \argmin_{t \in \R} Y(t) \label{eq: MedianDef}
\end{align}
als (irgendeine) Minimalstelle der Funktion $Y$.

\begin{definition}
	Hierbei heißt $m \in \R$ \begriff{Median} von der Zufallsgröße $X$.
\end{definition}

Wir haben den \begriff{Korrespondenzsatz} genutzt, der besagt, dass das Bildmaß $\P \circ  X^{-1}$ eindeutig von der Verteilungsfunktion $F$ bestimmt ist. Das heißt also wir schreiben
$F(\dx)$ und meinen $\P \circ X^{-1} (\dx)$

\textbf{Notation:} $F(m-) \defeq F(m-0) \defeq \lim_{t \uparrow m} F(t)$. Der Ausdruck ist wohldefiniert für Verteilungsfunktionen, denn diese sind rechtsseitig stetig und haben linksseitige Limiten.


In folgendem kleinen Lemma wollen wir die Menge aller Mediane charakterisieren.

\begin{lemma}\label{lemma: median}
	Sei $X \sim F_X \defqe F$ integrierbar und $m \in \R$. Dann sind äquivalent:
	\begin{enumerate}[label=(\arabic*)]
		\item \label{it:medianHalf} $F(m-)\le \frac{1}{2} \le F(m)$
		\item \label{it:medianMin} $\E\sqbrackets{\abs{X-t}}\ge\E{\abs{X-m}} \qquad\forall t \in \R$
		\item \label{it:medianMedian} $m$ ist Median
	\end{enumerate}
\end{lemma}
\begin{beispiel}
	Für $F$, $a < b \in \R$ mit
	\begin{align*}
		F(x) \begin{cases}
			< \frac12, &\text{ für } x < a \\
			= \frac12, &\text{ für } a \le x < b \\
			> \frac12, &\text{ für } x \ge b
		\end{cases}
	\end{align*}
	ist die Menge der Mediane genau das Intervall $[a, b]$.
	% Bild: Diagramm mit Funktion F mit Plateau bei 1/2 (1/2 eingetragen auf y-Achse
	% Punkt 1 exkl. unter 1/2, von Punkt 1 inkl. bis Punkt 2 exkl. Plateau,
	% an Punkt 2 Sprung, dann ansteigend
	% auf x-Achse Intervall [Punkt 1, Punkt 2] farbig = alle Mediane
\end{beispiel}

\begin{proof}
	\begin{description}
		\item[\ref{it:medianHalf} $\Rightarrow$ \ref{it:medianMin}:] Sei $Q \defeq \P \circ X^{-1}$ das zu $F$ gehörige Bildmaß. Setze
		\begin{equation*}
			h(t) \defeq \E\sqbrackets{\abs{X - t} - \abs{X - m}} \overset{lin.}{=} Y(t) - Y(m)
		\end{equation*}
		Dann ist \ref{it:medianMin} äquivalent zu $h(t) \ge 0$ für alle $t\in\R$. Dies wollen wir nun zeigen.
		\begin{enumerate}[label=Fall \Alph*:] 
			\item Sei $t < m$.
			\begin{align*}
					h(t) 
					\overset{\text{Trafo}}&{=}
					\int_{\R} \abs{x-t}-\abs{x-m} \, F(\dx)\\
					&=\int_{(-\infty,t]} \underbrace{\abs{x-t}-\abs{x-m}}_{= t-x-(m-x)=-(m-t)} \, F(\dx)
					+ \int_{(t,m)} \underbrace{\underbrace{\abs{x-t}-\abs{x-m}}_{\underbrace{x-t}_{\ge 0}-\underbrace{(m-x)}_{\le m-t}}}_{\ge-(m-t)} \, F(\dx)\\
					&\phantom{=}
					+\int_{[m,\infty)}\underbrace{\abs{x-t}-\abs{x-m}}_{x-t-(x-m)=m-t} \, F(\dx) \\
					&\ge-(m-t) * \underbrace{Q((-\infty,t])}_{F(t)}+
					\brackets{-(m-t) * F(m-)-F(t)}
					+(m-t) * \underbrace{Q([m,\infty))}_{1-\underbrace{Q((-\infty,m))}_{F(m-)}}\\
					&=-\underbrace{(m-t)}_{\ge 0} * (\underbrace{1-2 * F(m-)}_{\text{wegen 1.\ Ungl. in }\ref{it:medianHalf}: \ge 0})\\
					&\ge0
				\end{align*}
			\item Sei $t > m$.  Wir verfahren ganz ähnlich:
				\begin{align*}
					h(t)&=	\int_{(-\infty,m]} \abs{x-t} - \abs{x-m} F(\dx)
						+ \int_{(m,t]} \abs{x-t} - \abs{x-m} F(\dx) \\
					&\phantom{=}\ + \int_{(t,\infty)} \abs{x-t} - \abs{x-m} F(\dx) \\
					&= \int_{(-\infty,m]} t - x - (m-x) F(\dx)
						+ \int_{(m,t]} t - x - (x - m) F(\dx) \\
					&\phantom{=}\ + \int_{(t,\infty)} x - t - (x - m) F(\dx) \\
					&\ge (t-m)F(m) - (t-m)(F(t) - F(m)) + (m - t)(1 - F(t)) \\
					&= (t-m)(F(m) - F(t) + F(m) - 1 + F(t)) \\
					&=\underbrace{(t-m)}_{> 0}*(\underbrace{2 * F(m)-1}_{\text{wegen 2.\ Ungl.\ in }\ref{it:medianHalf}\ge0})\\
					&\ge 0
				\end{align*}
			\item Für $t = m$ ist die Aussage trivial.
		\end{enumerate}
		\item[\ref{it:medianMin} $\Rightarrow$ \ref{it:medianHalf}:]
			Nach Annahme ist $h(t) \ge 0$ für alle $t \in \R$.
			\begin{enumerate}[label=Fall \Alph*:] 
				\item Sei $t<m$. Die obige Rechnung im Fall 1 bei $\Rightarrow$ zeigt
				\begin{align*}
					0 \le h(t)
					&= -(m-t) * F(t)+ \int_t^m \underbrace{\underbrace{x}_{<m}-t-(m-x) }_{=2x-t-m \le m-t} F(\dx)+(m-t) * (1-F(m-)) \\
					& \le -(m-t)* \Big( F(t)-1 \underbrace{+F(m-)-F(m-)}_{=0}+F(t) \Big) \\
					&=\underbrace{(m-t)}_{>0}*(1-2* F(t)) \\
					&\Rightarrow\forall t<m: \enskip 0 \le(m-t)*(1-2* F(t))\\
					&\Rightarrow\forall t<m: \enskip 0 \le 1-2* F(t)\\
					&\Rightarrow\forall t<m: \enskip F(t) \le \frac{1}{2}\\
					&\overset{t \uparrow m}{\Rightarrow} F(m-)\le\frac{1}{2} \tag{Def. linksseitiger Limes}
				\end{align*}
			\item Sei $t > m$. Siehe 2. Fall, analog:
			\begin{align*}
				0 \le h(t)
				&= (t-m)* F(m) + \int_m^t \underbrace{t - x - (x - m) }_{=t + m - 2x \le t-m} F(\dx) + (m-t)*(1-F(t)) \\
				&\le (t-m)*\brackets{F(m)+ F(t) - F(m)-1 + F(t)} \\
				&= \underbrace{(t-m)}_{>0}(2F(t) - 1) \\
				&\Rightarrow \forall t<m: \enskip 0\le 2F(t) - 1\\
				&\Rightarrow \forall t<m: \enskip F(t) \ge \frac{1}{2}\\
				&\overset{t\downarrow m}{\Rightarrow} F(m) \ge \frac{1}{2} \tag{Rechtsstetigkeit von $F$}
			\end{align*}
			\end{enumerate}
		\item[\ref{it:medianMin} $\Rightarrow$ \ref{it:medianMedian}:]
			Die Aussage \ref{it:medianMin} ist offensichtlich äquivalent zur Definition des Medians.
	\end{description}
\end{proof}

\begin{bemerkung}~
	\begin{enumerate}[label=(\arabic*), leftmargin=*]
		\item \cref{lemma: median} besagt, dass $\menge{ m \in \R : m \text{ erfüllt \cref{it:medianHalf}}}$ die Menge aller Mediane von $F$ ist. Der Median ist im Allgemeinen nicht eindeutig bestimmt.
		\item Im Allgemeinen gibt es mehrere Mediane. Üblicherweise \textit{wählt} man $m \defeq F^{-1} \brackets{\frac{1}{2}}$, wobei
			\begin{align*}
				F^{-1}(u) \defeq \inf\menge{ x \in \R: F(x) \ge u} \quad \forall u \in (0,1)
			\end{align*}
			die \begriff{Quantilfuntion} oder auch \begriff{verallgemeinerte Inverse} ist. Für weiterführende Literatur siehe \cite[Seite 20]{witting1985mathStat}. Da
			\begin{align*}
				F\brackets{F^{-1}(u)-} \le u \le F\brackets{F^{-1}(u)} \qquad \forall u \in (0,1),
			\end{align*}
			erfüllt $m = F^{-1}\brackets{\frac{1}{2}}$ die Bedingung \ref{it:medianHalf} in \cref{lemma: median} und ist somit ein Median, nämlich der kleinste.
		\item Die obige Funktion \eqref{eq: DefY}
			\begin{align*}
				Y \colon \R \to \R \text{ mit } Y(t) = \int_\R \abs{x-t}~F(\dx) \qquad \forall t\in\R
			\end{align*}
			ist stetig\footnote{zur Stetigkeit von $Y$: nutze Folgenkriterium + dominierte Konvergenz mit Majorante $\abs{X} + \abs{t}$}, aber im Allgemeinen nicht differenzierbar, z.\,B.\ falls $F \sim X$ eine diskrete Zufallsvariable ist. In diesem Fall ist somit die Minimierung über Differentiation nicht möglich.
		\item Sei $\mu = \E(X)$ der Erwartungswert von $X$. Dann gilt (Übung):
			\begin{equation*}
				\mu = \argmin_{t \in \R} \E\sqbrackets{(X-t)^2} = \argmin_{t \in \R} \E\sqbrackets{(X-t)^2 - X^2}.
			\end{equation*}
			(Das zweite $X^2$ wird abgezogen, da das $\argmin$ nicht davon betroffen ist und so die Bedingung $\E\sqbrackets{X^2}<\infty$ entfällt.)
			Begründung:
			\begin{align*}
				\E\sqbrackets{(X-t)^2) - X^2}
				&= \E\sqbrackets{X^2 - 2tX + t^2 - X^2} \\
				&= \E\sqbrackets{X^2} - 2t\mu + t^2 - \E\sqbrackets{X^2}
				= t^2 - 2t\mu
			\end{align*}
			Minimiere nun diese quadratische Funktion und erhalte die gewünschte Resultat.
	\end{enumerate}
\end{bemerkung}

%\begin{notation}
	In der Statistik identifizieren wir oft stillschweigend Zufallsgrößen mit ihren Realisationen, also $X \leftrightsquigarrow x$ was sich formal $X(\omega_0) = x$ schreiben lässt.
%\end{notation}

Zur Schätzung von $m$ seien $X_1,\dots, X_n \sim F$ iid Zufallsvariablen mit zugehöriger \begriff{empirischer Verteilungsfunktion}
\begin{equation*}
	F_n \colon \R \to [0,1] 
	\quad \text{ mit } \quad
	F_n(x) \defeq \frac{1}{n} * \sum_{i=1}^n \one_{\menge{ x_i \le x}}(x) \qquad \forall x \in \R.
\end{equation*}
Tatsächlich ist $F_n$ die Verteilungsfunktion zum \begriff{empirischen Maß}
\begin{equation*}
	Q_n \colon \borel\R \to [0,1]
	\quad \text{ mit } \quad
	Q_n(B)\defeq\frac{1}{n}*\sum_{i=1}^n\delta_{x_i}(B)
\end{equation*}
wobei das \begriff{Dirac-Maß} in $t$ definiert ist als
\begin{equation*}
	\delta_t \colon \borel\R \to [0,1]
	\quad \text{ mit } \quad
	\delta_t(B) \defeq
	\begin{cases}
		0, &\falls t \notin B \\
		1, &\falls t \in B
	\end{cases}
	\qquad\forall B \in \borel\R \enskip \forall t \in \R
\end{equation*}

Gemäß dem Satz von Gliwenko-Cantelli gilt:
\begin{equation*}
	\sup_{x \in \R} \abs{F_n(x)-F(x)} \overset{n \to \infty}{\longrightarrow} 0 \qquad \P\text{-fast sicher für alle Verteilungsfunktionen } F
\end{equation*}
Die empirische Verteilungsfunktion konvergiert also gegen die wahre Verteilungsfunktion.

\begin{*erinnerung}
	Für das Dirac-Maß $\delta_x \colon \A \to \R_+$ mit $\delta_x(A) \defeq \one_A(x)$ gilt:
	\begin{equation}\label{eq: dirac}
		\int_\R f(t) \, \delta_x(\diff{t}) = f(x) \qquad \forall x \in \R
	\end{equation}
\end{*erinnerung}

\begin{*notation}
	Wir identifizieren eine Verteilung $\P\circ X^{-1}$ mit der zugehörigen Verteilungsfunktion $F_X$, also $\P \circ X^{-1} \leftrightarrow F_X$ und $F_X(\dx) \defeq (\P\circ X^{-1} \, \dx)$.
\end{*notation}

\begin{*erinnerung}
	Das Lebesgue-Maß ist linear im Maß, d.\,h.:
	\begin{equation}\label{eqLinMasse}
		\int_\omega f \diff{(a*\mu + b*\nu)} = a*\int_\omega f \diff{\mu} + b*\int_\omega f \, \diff{\nu}
	\end{equation}
\end{*erinnerung}
%
%Ein vages Stetigkeitsargument motiviert folgenden Schätzer für $m$,
%den wir durch Ersetzen von $F$ in $Y$ durch $F_n$ erhalte (\define{plug-in}-Methode):
%\begin{align*}
%	\hat{m}_n&\defeq\argmin_{t\in\R}Y_n(t)\defeq\text{ (irgendeine) Minimalstelle der Funktion}\\
%	Y_n(t)&\defeq\int_\R \abs{x-t}F_n(\d x)\\
%	&=\int_\R \abs{x-t}Q_n(\d x)\\
%	\overset{\Def~Q+\eqref{eqLinMasse}}&=
%	\frac{1}{n}*\sum_{i=1}^n\int_\R \abs{x-t}~\delta_{X_i}(\d x)\\
%	\overset{\eqref{eq: dirac}}&=
%	\frac{1}{n}*\sum_{i=1}^n\abs{X_i-t}
%\end{align*}
%
%$\hat{m}_n$ heißt \define{empirischer Median} von $X_1,\dots,X_n$ mit üblicher Auswahl $\hat{m}_n=F_n^{-1}\klammern{\frac{1}{2}}$ gemäß Lemma \ref{lemma: median} (da empirische Verteilungsfunktion eine Verteilungsfunktion ist).
%
%\begin{bemerkung}\
%	\begin{itemize}
%		\item Wenn man eine ungerade Anzahl von Daten hat, ist der Median der mittlere Wert, nachdem man die Daten der Größe nach geordnet hat.
%		\item Hat man hingegen eine gerade Anzahl an Daten, dann ist der Median der kleinere der beiden mittleren Werte.
%	\end{itemize}
%\end{bemerkung}
%
%Mit dem starken \undefine{Gesetz der großen Zahlen (SGGZ)} gilt
%\begin{align}\label{eq1.1}\tag{1.1}
%	\forall t \in \R: \klammern{Y_n(t)\ntoinf \Earg{\abs{X_1 - t}} = \int_{\R} \abs{x-t} F(\d x) = Y(t) \text{ $\P$-fast sicher}}
%\end{align}
%Problem: Folgt aus \eqref{eq1.1} bereits, dass
%\begin{align*}
%	\argmin_{t\in\R} Y_n(t)
%	\ntoinf
%	% CHECKED: '\longrightarrow' used.
%	\argmin_{t\in\R}
%	Y(t)\text{ fast sicher?}
%\end{align*}
%Wenn ja, dann hätten wir:
%\begin{align*}
%	\hat{m}_n
%	\ntoinf
%	% CHECKED: '\longrightarrow' used.
%	m\text{ fast sicher (\define{starke Konsistenz})}
%\end{align*}
%(Hierbei stellt sich auch wieder die Frage, was wir genau meinen, wenn $m$ nicht eindeutig ist.)
%
%Wir formalisieren und verallgemeinern:
%\begin{align*}
%	&X_i:(\omega, \A,P)\to(\R,\B(\R))\text{ messbar},\qquad\omega↦ X_i(\omega)\\
%	&\Rightarrow
%	Y_n(t)\defeqY_n(t,\omega)=
%	\frac{1}{n}*\sum_{i=1}^n\abs{X_i(\omega)-t}\\
%	&\Rightarrow
%	Y_n(t,*):(\omega,\A)\to(\R,\B(\R))\text{ messbar }\forall t\in\R
%\end{align*}
%
%\begin{defi}
%	Die \define{Kollektion}
%	\begin{align*}
%		Y_n\defeq\set{ Y_n(t,*):t\in\R}
%		=\set{ Y_n(t):t\in\R}
%	\end{align*}
%	heißt \define{stochastischer Prozess (SP)}. (Ein stochastischer Prozess
%	ist etwas allgemeiner eine Kollektion von Zufallsvariablen, die von
%	einem Parameter abhängen, hier $t\in\R$.) Die Abbildung
%	\begin{align*}
%		Y_n(*,\omega):\R\to\R,\qquad t↦ Y_n(t,\omega)
%	\end{align*}
%	heißt \define{Trajektorie/ Pfad} des SP $Y_n$ zu festem $\omega\in\omega$.
%\end{defi}
%
%In unserem Beispiel sind für \emph{alle} $\omega\in\omega$ die Pfade stetig auf $\R$. Die Abbildung
%\begin{align*}
%	Y_n:\omega\to C(\R;\R),\qquad\omega↦ Y_n(*,\omega)
%\end{align*}
%heißt \define{Pfadabbildung} des SP $Y_n$. Wir identifizieren also den SP $Y_n$ mit seiner Pfadabbildung. Damit ist $Y_n$ eine Abbildung von $\omega$ in den Funktionenraum
%\begin{align*}
%	C(\R)\defeqC(\R;\R)\defeq\set{ f:\R\to\R: f\text{ ist stetig}}.
%\end{align*}
%
%Sei $d:C(\R)× C(\R) → \R$ die Metrik der gleichmäßigen Konvergenz auf Kompakta auf $C(\R)$ (formale Definition kommt später) und sei
%\begin{align*}
%	\B(C(\R))
%	\defeq\B_d\klammern[\big]{C(\R)}
%	% CHECKED: '\big' used.
%	\defeq σ \klammern[\big]{\set{ G⊆ C(\R):G\text{ ist offen bzgl. }d}}
%	% CHECKED: '\big' used.
%\end{align*}
%die von $d$ induzierte \define{Borel-$σ$-Algebra}.
%
%Wir werden sehen, dass die Abbildung
%\begin{align*}
%	Y_n:(\omega,\A,\P)\to\klammern[\Big]{C(\R),\B\klammern[\big]{C(\R)}}
%	% CHECKED: '\big' '\Big' used.
%\end{align*}
%messbar ist. $Y_n$ ist also eine Zufallsvariable mit Werten im metrischen Raum $(C(\R),d)$.
%
%\paragraph{Formulierung des Problems im allgemeinen Rahmen}
%Seien $Y_n,~n\inℕ$ mit $Y$ SP mit stetigen Pfaden (\define{stetige SP}).
%Was lässt sich sagen über die Gültigkeit der folgenden Implikation?
%\begin{align}\label{eq1.2}\tag{1.2}
%	Y_n
%	\ntoinf
%	% CHECKED: '\longrightarrow' used.
%	Y\text{ fast sicher }
%	\Rightarrow
%	\argmin_{t\in\R} Y_n(t)
%	\ntoinf
%	% CHECKED: '\longrightarrow' used.
%	\argmin_{t\in\R} Y(t)\text{ fast sicher}
%\end{align}
%\subparagraph{Ziel} Welche Art der Konvergenz $Y_n\ntoinf  Y$ reicht für obige Implikation aus? Gleichmäßige Konvergenz, gleichmäßige Konvergenz auf Kompakta, punktweise Konvergenz oder sogar nur \eqref{eq1.1}?
%% CHECKED: '\longrightarrow' used.
%
%$Y$ besitzt womöglich (mit positiver Wahrscheinlichkeit) keine eindeutige Minimalstelle. Und dann?
%
%Betrachten wir nun wieder eine andere Art der Konvergenz:
%Für die Konstruktion von (asymptotischen) Konfidenzintervallen für $m$ benötigt man (in der Statistik) \define{Verteilungskonvergenz} (Wiederholung: $\distrto $ ist die Notation für \undefine{konvergiert in Verteilung}):
%\begin{align}\label{eq1.3}\tag{1.3}
%	a_n(\hat{m}_n-m)
%	\distrto ζ\text{ in }\R
%	% CHECKED: '\longrightarrow' used.
%\end{align}
%wobei $a_n\to\infty$ eine normalisierende Folge ist und $ζ$ eine Grenzvariable ist, die es zu identifizieren gilt.
%
%\begin{beispiel}[Zentraler Grenzwertsatz für den Durchschnitt]
%	Für den Mittelwert von $\overline{X_n}$ von integrierbaren ZV $(X_n)_n$ mit Mittelwert $\mu$ ist diese Folge $a_n = √{n}$. Dafür besagt das Gesetz der großen Zahlen, dass
%	\begin{equation*}
%		\overline{X_n} - \mu → 0 \text{ fast sicher}.
%	\end{equation*}
%	Der \undefine{zentrale Grenzwertsatz} liefert uns
%	\begin{equation*}
%		√{n}\klammern{\overline{X_n} - \mu} \overset{\L}{→} N(0, σ^2)
%	\end{equation*}
%	also die Konvergenz in Verteilung gegen die Normalverteilung.
%\end{beispiel}
%Für die Herleitung von \eqref{eq1.3} favorisiere wieder einen \undefine{funktionalen Ansatz}. Sei
%\begin{align*}
%	Z_n(t)\defeqβ_n*\klammern{ Y_n\klammern{m+\frac{t}{a_n}}-Y_n(m)}\qquad t\in\R
%\end{align*}
%der sogenannte \define{reskalierte Prozess zu $Y_n$}, wobei $β_n$ (zunächst) irgendeine positive Folge ist. Damit folgt
%\begin{align}\label{eq1.4}\tag{1.4}
%	σ_n \defeq a_n(\hat{m}_n-m) \in \argmin_{t\in\R} Z_n(t),
%\end{align}
%denn:
%\begin{proof}
%  Zu zeigen ist:
%	\begin{align*}
%		&& Z_n(σ_n) &\le Z_n(t) & \forall t \in \R \\
%		&⇔& β_n \klammern{Y_n\klammern{m + \frac{σ_n}{a_n}} - Y_n(m)} &\le Z_n(t) & \forall t \in \R\\
%		&⇔& β_n \klammern{Y_n\klammern{m + \hat{m}_n - m} - Y_n(m)}
%		&\le β_n \klammern{Y_n\klammern{ m + \frac {t}{a_n}}- Y_n(m)} & \forall t \in \R \\
%		&⇔& β_n\klammern{Y_n(\hat m_n) - Y_n(m)}
%		&\le β_n \klammern{Y_n\klammern{ m + \frac {t}{a_n}} - Y_n(m)} & \forall t \in \R \\
%		&⇔& Y_n(\hat m_n) &\le Y_n\klammern{m + \frac t {a_n}} & \forall t \in \R \\
%		&⇔& Y_n(\hat m_n) &\le Y_n(u) & \forall u \in \R.
%	\end{align*}
%	Das gilt aber, weil $\hat m_n \in \argmin_{t \in \R} Y_n(t)$.
%\end{proof}
%Klar: $Z_n$ ist wieder ein stetiger stochastischer Prozess und damit $(Z_n)_{n\inℕ}$ eine Folge von Zufallsvariablen in $(C(\R),d)$.
%% CHECKED: '\big' used.
%Wünschenswert auch hier wäre die Gültigkeit folgender Implikation:
%\begin{align}\label{eq1.5}\tag{1.5}
%	% CHECKED: ist Tatsächlich das L hier bei den argmin gemeint? Macht das Sinn, sind die argmin auch eine Folge von Zufallsvariablen? -> ja, sind sie
%	Z_n
%	\distrto
%	% CHECKED: '\longrightarrow' used.
%	Z\text{ in } (C(\R),d)
%	\Rightarrow
%	\argmin_{t\in\R} Z_n(t)
%	\distrto
%	% CHECKED: '\longrightarrow' used.
%	\argmin_{t\in\R} Z(t)
%\end{align}
%Da wir $σ_n = a_n(\hat m_n - m) → ζ$ zeigen wollen und $\argmin_{t\in\R}Z_n(t) = σ_n$, haben wir jetzt die erste Charakterisierung von $ζ$ gefunden.
%
%Für die Aussage \ref{eq1.5} ist erforderlich, das Konzept der \define{Verteilungskonvergenz von Zufallsvariablen in metrischen Räumen}, damit \eqref{eq1.5} eine wohldefinierte Bedeutung erhält.
%Dies folgt später.
%
%Natürlich auch hier wieder das Problem: $Z$ besitzt mit positiver Wahrscheinlichkeit mindestens 2 Minimalstellen.
%Und dann?
%
%Im Falle einer fast sicher eindeutigen Minimalstelle von $Z$ würde aber aus \eqref{eq1.4} und \eqref{eq1.5} folgen:
%\begin{align*}
%	a_n(\hat{m}_n-m)
%	\distrto
%	% CHECKED: '\longrightarrow' used.
%	\argmin_{t\in\R} Z(t)
%\end{align*}

    
    \chapter{Metrische Räume}
    \label{chapter_2_metrischeRaeume}
    
Sei $(\S,d)$ ein metrischer Raum.
%
\begin{beispiel}[Supremums-Metrik] %2.1
	\begin{align*}
		\S = C([0,1]) & \defeq \menge{f \colon [0,1] \to \R : f\text{ stetig}} \\
		d(f,g) &\defeq \sup_{t \in [0,1]} \abs{f(t)-g(t)} \qquad \forall f,g \in C([0,1])
	\end{align*}
\end{beispiel}

\begin{definition}
	\begin{enumerate}[label={(\arabic*)}]
		\item Für $x \in\S$ und $r>0$ ist
		\begin{align*}
			B(x,r) \defeq B_d(x,r) \defeq \menge{y \in \S : d(x,y) < r}
		\end{align*}
		die \begriff{offene Kugel um Mittelpunkt $x$ und Radius $r$}.
		\item Sei $A \subseteq \S$. Dann ist:
		\begin{itemize}[nolistsep]
			\item $\inner A \dots$ das \begriff{Innere} von $A$
			\item $\abschluss{A} \dots$ die \begriff{abgeschlossene Hülle} von $A$
			\item $\rand A \defeq \abschluss{A} \cap \abschluss{A^C} =  \abschluss{A} \setminus \inner A \dots$ der \begriff{Rand} von $A$
			\item $A^\complement \defeq \S \setminus A \dots$ das \begriff{Komplement} von $A$
		\end{itemize}
		\item Die durch $d$ induzierte \begriff{Topologie} ist
		\begin{align*}
			\mathcal{G} \defeq \mathcal{G}(\S) &\defeq \menge{G \subseteq \S: G \text{ ist offen bzgl. } d} \\
			&\phantom{:}= \menge{G \subseteq \S: \forall x \in G: \exists \, r > 0 : B_d(x,r) \subseteq G}
			%
			\intertext{und}
			%
			\mathcal{F} \defeq \mathcal{F}(\S) 
			&\defeq \menge{F \subseteq \S : F \text{ ist abgeschlossen}} 
			= \menge{F : F = G^\complement, G \in \mathcal{G}}
		\end{align*}
		ist die \begriff{Menge der abgeschlossenen Mengen}.
		\item Sei $\emptyset \neq A \subseteq \S$ und $x \in \S$. Dann ist $d(x,A) \defeq \inf\menge{d(x,a) : a \in A} \ge 0$ der \begriff{Abstand von $x$ zu $A$}.
		\item Es ist
		\begin{align*}
			C(\S) &\defeq \menge{f \colon S \to \R : f \text{ stetig}} \\
			C^b(\S) &\defeq \menge{f \in C(\S) : f \text{ beschränkt}} \\
			\norm{f} &\defeq \norm{ f}_\infty \defeq \sup_{x\in\S} \abs{f(x)}
		\end{align*}
	\end{enumerate}
\end{definition}

\begin{lemma} \label{lemma: 2.3} Sei $A \neq \emptyset$.
	\begin{enumerate}[label={(\arabic*)}]
		\item \label{it: distCharakterisierung} Es ist $x \in\abschluss{A} \equivalent d(x,A)=0$.
		\item \label{it: distDreieck} Für alle $x,y \in \S$ gilt $\abs{d(x,A)-d(y,A)} \le d(x,y)$.
		\item \label{it: distStetig} Die Abbildung $d(*, A) \colon \S \to [0,\infty)$, $x \mapsto d(x,A)$ ist gleichmäßig stetig.
	\end{enumerate}
\end{lemma}

\begin{proof}
	\begin{enumerate}[label=(zu \arabic*), leftmargin=*]
		\item \begin{equivalence}
			\hinrichtung Sei $x \in \abschluss{A}$. Dann existiert für alle $\epsilon > 0$ ein $a \in A$ mit $d(x,a) < \epsilon$. Somit 
			\begin{equation*}
				0 \le d(x,A) \le d(x,a) < \epsilon \quad \forall \epsilon > 0 
				\overset{\epsilon \to 0}{\follows} d(x,A) = 0
			\end{equation*}
			\rueckrichtung Sei $d(x,A)=0$. Dann folgt aus der Infimumseigenschaft:
			\begin{equation*}
				\forall \epsilon > 0 \enskip \exists \, a \in A : \quad 0 \le d(x,a) \le 0 + \epsilon = \epsilon \\
				\follows x \in \abschluss{A}
			\end{equation*}
		\end{equivalence}
		\item Seien $x,y \in \S$. Dann gilt $d(x,a) \le d(x,y) + d(y,a)$ mit der Dreiecksungleichung für alle $a \in A$ und somit
		\begin{equation*}
			d(x,A) \le d(x,y) + d(y,A) 
			\follows d(x,A) - d(y,A) \le d(x,y)
		\end{equation*}
		Vertauschen von $x$ und $y$ liefert $d(y,A) - d(x,A) \le d(y,x) = d(x,y)$ und daraus folgt die Behauptung.
		\item Folgt aus \cref{it: distDreieck} da die Funktion $d(\cdot,A)$ Lipschitz-stetig und damit gleichmäßig stetig ist.
	\end{enumerate}
\end{proof}

\begin{satz}[Stetige Approximation von Indikatorfunktionen] \label{satz: 2.4}
	Zu $A \subseteq \S$ und $\epsilon > 0$ existiert eine gleichmäßig stetige Funktion
	\begin{equation*}
		f \colon \S \to [0,1] \quad \text{ mit der Eigenschaft} \quad f(x) =
		\begin{cases}
			1 & \falls x \in A \\
			0 & \falls d(x,A) \ge \epsilon
		\end{cases}
	\end{equation*}
\end{satz}
\begin{proof}
	Setze
	\begin{equation*}
		\phi \colon \R \to [0,1] 
		\quad \mit \quad 
		\phi(t) \defeq
		\begin{cases}
			1    & \falls t \le 0 \\
			1-t  & \falls 0 < t < 1 \\
			0    & \falls t \ge 1
		\end{cases}
	\end{equation*}
	Dann ist $\phi$ gleichmäßig stetig auf $\R$. Sei
	\begin{equation*}
		f(x) \defeq \phi \brackets{\frac{1}{\epsilon} *  d(x,A)} \qquad \forall x \in \S
	\end{equation*}
	Dann hat dieses $f$ die gewünschte Eigenschaft wegen \cref{lemma: 2.3}.
\end{proof}
%
\begin{definition}\label{def: 2.5}
	Ein metrischer Raum $(\S,d)$ heißt \begriff{separabel}
	\begin{align*}
		&:\Leftrightarrow \exists \text{ abzählbares } S_0 \subseteq \S :\S \subseteq \abschluss{S_0} \\
		&\phantom{:} \Leftrightarrow \exists \text{ abzählbares } S_0 \subseteq \S : \S = \abschluss{S_0} \\
		&\phantom{:}\Leftrightarrow \exists \text{ abzählbares } S_0 \subseteq \S : S_0 \text{ liegt dicht in } \S
	\end{align*}
\end{definition}

\begin{beispiel} \label{bsp: 2.6}
	Der metrische Raum $C([0,1])$ mit Supremums-Metrik ist separabel.
\end{beispiel}
\begin{proof}
	Die Menge $S_0 \defeq \menge{P : P \text{ ist Polynom mit rationalen Koeffizienten}}$ ist abzählbar. Aus dem \textit{Approximationssatz von Weierstraß} und der Dichtheit von $\Q$ folgt die Behauptung.
\end{proof}

\begin{definition} %2.7
	$\G_0 \subseteq \G$ heißt \begriff{Basis} von $\G$ genau dann, wenn sich jedes $G \in \G$ als Vereinigung von Mengen aus $\G_0$,
	so genannte \begriff{$\G_0$-Mengen}, darstellen lässt.
\end{definition}
%
\begin{beispiel}\label{bsp: 2.8} %2.8
	Die Menge $\menge{ B(x,r) : x \in \S, 0 < r \in \Q}$ ist Basis von $\G$.
\end{beispiel}
\begin{proof}
	Sei $G \in \G$. Dann gilt:
	\begin{equation*}
		\forall x \in G \enskip \exists \,  0 < r_x \in \Q : B(x,r_x) \subseteq G 
		\follows
		G = \bigcup_{x\in G} \menge{x}
		\subseteq \bigcup_{x \in G} \underbrace{B(x,r_x)}_{\subseteq G} \subseteq G 
%		\follows 
%		G = \bigcup_{x\in G} \underbrace{B(x,r_x)}_{\in \G_0}
	\end{equation*}
	Damit gilt also Gleichheit.
\end{proof}


\begin{satz} \label{satz: 2.9}
	$\S$ separabel $\Leftrightarrow$ $\G$ hat eine abzählbare Basis
\end{satz}

\begin{proof}
	\begin{equivalence}
		\hinrichtung Sei $S_0\subseteq \S$ abzählbar und dicht in $\S$. Wir zeigen, dass
		\begin{equation*}
			\G_0 \defeq \menge{ B(x,r) : x \in S_0 , 0 < r \in \Q}\subseteq \G
		\end{equation*}
		eine Basis ist.
		Sei also $G$ offen. Dann folgt aus Beispiel \cref{bsp: 2.8}, dass
		\begin{equation} \label{proof: 2.9-sternchen} \tag{$\star$}
			G = \bigcup_{x \in G} B(x,r_x), \qquad 0 < r_x \in\Q \enskip \forall x \in G
		\end{equation}
		Da $\quer{S_0} = \S$ gilt:
		\begin{align*}
			&\forall x \in G:\exists y_x \in S_0: d(x,y_x) < \frac{r_x}{2} \\
			&\Rightarrow \quad d(x,y)
			\overset{\triangle-\text{Ungl.}}{\le}
			d(x,y_x) + d(y_x,x) < \frac{r_x}{2} + \frac{r_x}{2} = r_x
			\qquad \forall y \in B\brackets{y_x,\frac{r_x}{2}} \\
			&\Rightarrow \quad B\brackets{y_x,\frac{r_x}{2}} \subseteq B(x,r_x)
			\qquad \forall x \in G \\
			&\Rightarrow \quad G \overset{\eqref{proof: 2.9-sternchen}}{\supseteq}
			\bigcup_{x \in G} \underbrace{ B\brackets{y_x,\frac{r_x}{2}} }_{\supseteq \menge{x}}
			\supseteq \bigcup_{x \in G} \menge{x} = G % \\
%			&\Rightarrow \quad G = \bigcup_{x\in G} \underbrace{B\brackets{y_x,\frac{r_x}{2}}}_{\in \G_0}
		\end{align*}
		Damit gilt also Gleichheit und $\G_0$ ist eine Basis. Da $S_0$ abzählbar ist, ist $\G_0$ abzählbar. 
		%
		\rueckrichtung Sei $\G_0$ abzählbare Basis von $\G$ und sei oBdA $\emptyset \notin \G_0$. Wähle für jedes $G \in \G_0$ ein $x_G \in G$ fest aus. Setze
		$S_0 \defeq \menge{x_G : G \in \G_0}$. Dann ist $S_0$ auch abzählbar. Es verbleibt die Dichtheit zu zeigen.
		Sei $x \in \S$ und $\epsilon > 0$ beliebig. Da $B(x,\epsilon)$ offen und $\G_0$ eine Basis ist, gilt:
		\begin{equation*}
			\exists \, \G_{x,\epsilon}\subseteq\G_0\mit B(x,\epsilon)=\bigcup_{G\in\G_{x,\epsilon}} G
			\follows  G\subseteq B(x,\epsilon)\qquad\forall G\in\G_{x,\epsilon}
		\end{equation*}
		Wähle ein $G$ von diesen aus. Dann gilt:
		\begin{equation*}
			x_G \in G \subseteq B(x,\epsilon)
			\follows  x_G\in B(x,\epsilon)
			\follows  d(\underbrace{x_G}_{\in S_0},x) < \epsilon
		\end{equation*}
	\end{equivalence}	
\end{proof}

\begin{satz}\label{satz: 2.10} %2.10
	Seien $(\S,d)$ und $(\S',d')$ metrische Räume.
	\begin{enumerate}[label={(\arabic*)}]
		\item \label{it: produktmetrik} Auf $\S \times \S'$ sind Metriken definiert durch
			\begin{align*}
				d_1 \brackets{(x,x'),(y,y')}
				&\defeq \brackets{ \brackets{d(x,y)}^2 + \brackets{d'(x',y')}^2}^{\frac{1}{2}}
				&\forall(x,x'),(y,y') \in \S \times \S'\\
				d_2 \brackets{(x,x'),(y,y')}
				&\defeq \max\menge{ d(x,y),d'(x',y')}
				&\forall(x,x'),(y,y')\in \S \times \S' \\
				d_3\brackets{(x,x'),(y,y')}
				&\defeq d(x,y) + d'(x',y')
				&\forall(x,x'),(y,y') \in \S \times \S'
			\end{align*}
		\item \label{it: produkttopologie} Die Metriken $d_1$, $d_2$ und $d_3$ induzieren dieselbe Topologie $\mathcal{G}(\S \times \S')$ auf $\S\times\S'$,
			die sogenannte \begriff{Produkttopologie} von $\mathcal{G}(\S)$ und $\mathcal{G}(\S')$.
		\item \label{it: produktbasis}
			$\mathcal{G}(\S\times\S')=\menge{\bigcup_{
					\begin{subarray}{c} G\in\mathcal{O}\\ G'\in\mathcal{O}' \end{subarray}
				} G\times G': \mathcal{O}\subseteq\mathcal{G}(\S), \mathcal{O}'\subseteq\mathcal{G}(\S')
				}$, d.h.
			\begin{equation}
				\menge{ G \times G' : G \in \mathcal{G}(\S), G' \in \mathcal{G}(\S')}
				\tag{Menge offener Rechtecke}
			\end{equation}
			bildet eine Basis von $\mathcal{G}(S\times\S')$.
	\end{enumerate}
\end{satz}

\begin{proof}
	\begin{enumerate}[label=(zu \arabic*), leftmargin=*]
		\item \label{it: zu_ref_it_produktmetrik}
		Überprüfung der Eigenschaften einer Metrik (zur Übung).
%		
		\item \label{par:zu_ref_it_produkttopologie}
			Punktweise gelten die Beziehungen:
			\begin{equation*}
				d_2 \le d_1 \le\sqrt{2}* d_2,
				\qquad
				\frac{1}{\sqrt{2}} * d_3 \le d_1 \le d_3,
				\qquad
				d_2 \le d_3 \le 2 * d_2
			\end{equation*}
			Beachte beim Nachweis, dass die $d_i$'s als Metriken größer oder gleich Null sind. Aus obigen Beziehungen folgt u.\,a.:
			\begin{equation*}
				B_{d_2}\brackets{x,\frac{r}{\sqrt{2}}} \subseteq  B_{d_1}(x,r)
			\end{equation*}
			denn
			\begin{equation}
				r > \sqrt{2} * d_2(y,x) \ge d_1(y,x) \label{eq: metrikvergleich}
			\end{equation}
			Damit folgt aus $G$ $d_1$-offen schon, dass $G$ $d_2$-offen ist, denn für $x \in G$ existiert $r > 0$ mit $B_{d_1}(x, r) \subseteq G$ und \eqref{eq: metrikvergleich} liefert
			$B_{d_2}(x, \frac{r}{\sqrt{2}}) \subseteq G$. Damit ist $x$ auch ein ein $d_2$-innerer Punkt. Die anderen Relationen gelten analog.
			%
		\item \begin{description}
				\item[$\subseteq$:] 
				\label{it: zu_ref_it_produktbasis}
				Sei $G^\ast \in \mathcal{G}(\S \times \S')$ eine offene Menge in der Produkttopologie. Dann gilt
				\begin{equation*}
					\forall x^\ast = (x,x') \in G^\ast: \, \exists r = r_{x^\ast} > 0: \enskip 
					G^\ast = \bigcup_{x^\ast \in G^\ast} B\brackets{x^\ast,r_{x^\ast}}.
				\end{equation*}
				Wegen \cref{it: produkttopologie} sei oBdA. $\S^\ast \defeq \S \times \S'$ versehen mit der Metrik $d_2$. Dann gilt:
				\begin{align*}
					B_{d_2}\brackets{x^\ast,r_{x^\ast}}
					&= \menge{(y,y')\in \S \times \S': \max \menge{d(x,y),d'(x',y')} < r_{x^\ast}} \\
					&= \menge{(y,y') \in \S \times \S': d(x,y) < r_{x^\ast} \land d'(x',y') < r_{x^\ast}} \\ 
					&=  \underbrace{B_d\brackets{x,r_{x^\ast}}}_{\in \mathcal{G}(\S)} \times \underbrace{B_{d'}\brackets{x', r_{x^\ast}}}_{\in\mathcal{G}(\S')}
				\end{align*}
				\item[$\supseteq$:] Sei zunächst $G \times G'$ $G,G'$ offen und $x^\ast = (x,x') \in G \times G'$. Also ist $x \in G$ und $x' \in G'$ und somit
				\begin{equation*}
					\exists \, r,r' > 0: B_d(x,r) \subseteq G \land B_{d'}(x',r') \subseteq G'
				\end{equation*}
				Setze $r^\ast \defeq \min\menge{ r,r'}>0$. Damit folgt
				\begin{align*}
					B_{d_2}\brackets{x^\ast,r^\ast}
					&\subseteq B_d(x,r) \times B_{d'}\brackets{x',r'} \\
					&\subseteq G \times G' = G^\ast \\
					&\Rightarrow G\times G' \in \mathcal{G}(\S \times \S')\\
					&\Rightarrow  \bigcup_{\begin{subarray}{c} G \in \mathcal{O} \\ G' \in \mathcal{O}' \end{subarray}} G \times G'\subseteq \mathcal{G}(\S\times\S')
					\qquad \forall \mathcal{O} \subseteq \mathcal{G}(\S), \mathcal{O}' \subseteq \mathcal{G}(\S')
				\end{align*}
				da die Produkttopologie vereinigungsstabil ist.
			\end{description}
	\end{enumerate}		
\end{proof}

\begin{definition}
	Die Metriken $d_1, d_2$ und $d_3$ heißen \begriff{Produktmetriken}.
	Daher alternative Schreibweise $d \times d'$, also z.\,B.\ $d\times d' \defeq \max\menge{d, d'}$ usw.
\end{definition}

\begin{bemerkung} %2.11
	Analog zu obigen Definitonen lassen sich Produktmetriken für endlich viele metrische Räume $(\S_i, d_i)_{i \in \menge{1, \dots ,k}}$ definieren, z.\,B.
	\begin{equation*}
		d_1 \times \dots \times d_k \defeq \brackets{\sum_{i=1}^k d_i^2}^{\frac{1}{2}},
	\end{equation*}
	die wiederum dieselbe Produkttopologie induzieren.
	Die bisherigen Resultate gelten analog.
\end{bemerkung}


	\chapter{Zufallsvariablen in metrischen Räumen}
	\label{chapter_3_zufallsvariablenInMetrischenRaeumen}
	% This work is licensed under the Creative Commons
% Attribution-NonCommercial-ShareAlike 4.0 International License. To view a copy
% of this license, visit http://creativecommons.org/licenses/by-nc-sa/4.0/ or
% send a letter to Creative Commons, PO Box 1866, Mountain View, CA 94042, USA.

\section{Messbarkeit in metrischen Räumen}
\label{sec: 3.1}

\begin{definition}\label{definition: 3.1}
	Die durch den metrischen Raum $(\S,d)$ induzierte $\sigma$-Algebra
	\begin{align*}
		\B \defeq \B(\S) \defeq \B_d(\S) \defeq \sigma \brackets{(\G(\S)} \defeq \sigma(\G).
	\end{align*}
	heißt \begriff{Borel-$\sigma$-Algebra}.
	Elemente $B  \in  \B(\S)$ heißen \begriff{Borel-Mengen} in $\S$.
\end{definition}

Beachte: $\B(\S) = \B_d(\S)$ hängt im Allgemeinen von der Metrik $d$ ab.

\begin{lemma}\label{lemma: 3.2}
	Es gilt:
	\begin{enumerate}[label=(\arabic*)]
		\item \label{it: 3.2-OpenClosed} 
		$\B(\S) = \sigma\brackets{\mathcal{F}(\S)}$
		\item \label{it: 3.2-StetigMessbar} Ist $f \colon (\S,d) \to (\S',d)$ stetig, so ist $f$ auch $\B_d(\S)$-$\B_d(\S')$-messbar.
		\item \label{it: 3.2-BasisSigmaErzeuger}
			Sei $\G_0$ abzählbare Basis von $\G(\S)$. Dann gilt $\sigma(\G_0) = \B(\S)$.
	\end{enumerate}
\end{lemma}

\begin{proof}
	\begin{enumerate}[label=(zu \alph*), leftmargin=*]
		\item \begin{itemize}
			\item[$\subseteq$:] Sei $G \in \G(\S)$. Dann gilt ist $G^\complement \in \mathcal{F}(\S) \subseteq \sigma\brackets{\mathcal{F}(\S)}$. Da $\sigma\brackets{\mathcal{F}(\S)}$ stabil unter Bildung von Komplementen ist, ist auch $G = \brackets{G^\complement}^\complement \in \sigma(\mathcal{F}(\S))$. Weiter folgt aus $\G \subseteq \sigma(\mathcal{F})$ schon $\sigma(\G) \subseteq \sigma(\mathcal{F})$.
			\item[$\supseteq$:] analog
		\end{itemize}
		\item Per Definition gilt $f^{-1}\brackets{\B_{d'}(\S')}
		= f^{-1}\brackets{\sigma\brackets{\G(\S')}}$. 
		Ein Satz aus der Maßtheorie liefert uns weiter $f^{-1}\brackets{\sigma\brackets{\G(\S')}} = \sigma\brackets{f^{-1}\brackets{\G(\S')}}$. 		
		Bekannterweise sind für stetige Funktionen Urbilder offener Mengen wieder offen, d.h. es gilt $f^{-1}\brackets{\G(\S')} \subseteq \G(S)$. 
		Somit ist $\sigma\brackets{f^{-1}\brackets{\G(\S')}} \subseteq \sigma\brackets{\G(\S)} = \borel{\S}$.
		\item \begin{itemize}
			\item[$\subseteq$:] klar wegen $\G_0 \subseteq \G(\S)$ und $\sigma$ monoton
			\item[$\supseteq$:] Sei $G \in \G$. Dann existieren geeignete $G_i \in \G_0 \subseteq \sigma(\G_0)$, sodass $G = \bigcup_{i \in \N} G_i$. Damit ist $G \in \sigma(\G_0)$.
			Aus der Stabilität unter Vereinigungen folgt die Behauptung:
			\begin{align*}
				G=\bigcup_{i \in \N} G_i \subseteq \sigma\brackets{\G_0}=\B(\S) \in \sigma(\G_0),
			\end{align*}
			da $\sigma$-Algebren \textit{abzählbar} vereinigungsstabil sind.
		\end{itemize}
	\end{enumerate} 
\end{proof}

\begin{satz}\label{satz: 3.3}
	Sei $(\S,d)$ separabler metrischer Raum. Dann gilt:
	\begin{align*}
		\B_{d \times  d}(\S \times \S) = \B(S) \otimes \B(\S)
	\end{align*}
	Ohne Separabilität gilt nur \enquote{$\supseteq$}.
	Für $S_1$ und $S_2$ separabel gilt $\B_{d_1  \times  d_2}(\S_1 \times  \S_2) = \B(S_1) \otimes \B(\S_2)$ analog, auch erweiterbar auf endliche Produkte.
\end{satz}

\begin{proof}
	Seien
	\begin{align*}
		\pi_1 &\colon \S \times \S \to \S \mit \pi_1(x,y) \defeq x \qquad \forall(x,y) \in \S \times \S \\
		\pi_2 &\colon \S \times \S \to \S \mit \pi_2(x,y) \defeq y \qquad \forall(x,y) \in \S \times \S 
	\end{align*}
	die \begriff{Projektionsabbildungen}. Dann gilt
	\begin{align*}
		\B(\S) \otimes \B(\S)
		&= \sigma(\pi_1,\pi_2) 
		= \sigma \brackets{\pi_1^{-1} \brackets{\sigma(\G)} \cup \pi_2^{-1} \brackets{\sigma(\G)}} \tag{Definition} \\
		&= \sigma\brackets{\sigma\brackets{\pi_1^{-1}(\G)} \cup \sigma\brackets{\pi_2^{-1}(\G)}} \\
		\overset{(\star)}&{=}
		\sigma\brackets{\pi_1^{-1}(\G) \cup \pi_2^{-1}(\G)} \\
		&= \sigma\brackets{\menge{ G \times  S,S \times  G' : G,G' \in \G}} \\
		&= \sigma\brackets{\menge{ \overbrace{G \times  G'}^{=(G \times  S) \cap (S \times  G')} : G,G' \in \G}} 
		&& \tag{\enquote{$\subseteq$}, da $S  \in  \G$ ; \enquote{$\supseteq$}, da $\sigma$-Algebra $\cap$-stabil} \\
		\overset{(\star\star)}&=
		\sigma\brackets{\menge{
		\bigcup_{
			\begin{subarray}{c}
				G \in \mathcal{O}\\
				G' \in \mathcal{O}'
			\end{subarray}
		}
		G \times  G' : \mathcal{O} , \mathcal{O}' \subseteq \G
		}} \\
		&=
		\sigma\brackets{\G(\S \times \S)} \tag{\cref{satz: 2.10}, \cref{it: produktbasis}} \\
		\overset{\text{Def}}&{=}
		\B(\S \times \S)
	\end{align*}
	
	Zum Nachweis von ($\star$):
	\begin{itemize}
		\item[$\supseteq$:] Setze $\mathcal{E} \defeq
		\underbrace{\sigma\brackets{\pi_1^{-1}(\G)}}_{\supseteq \pi_1^{-1}(\G)}
		\cup \underbrace{\sigma\brackets{\pi_2^{-1}(\G)}}_{\supseteq \pi_2^{-1}(\G)}
		\supseteq \pi_1^{-1}(\G) \cup \pi_2^{-1}(\G) \defqe \mathcal{H}$. Dann ist $\sigma(\mathcal{E}) \supseteq \sigma(\mathcal{H})$.
		\item[$\subseteq$:] Es ist $\pi_1^{-1}(\G) \subseteq \brackets{\pi_1^{-1}(\G) \cup \pi_2^{-1}(\G)} = \mathcal{H}$. Also ist auch $\sigma\brackets{\pi_1^{-1}(\G)} \subseteq \sigma(\mathcal{H})$. Analog erhalten wir auch $\sigma\brackets{\pi_2^{-1}(\G)} \subseteq \sigma(\mathcal{H})$. Dann gilt
		\begin{equation*}
			\mathcal{E} = 		\underbrace{\sigma\brackets{\pi_1^{-1}(\G)}}_{ \subseteq \sigma(\mathcal{H})}
			 \cup \underbrace{\sigma\brackets{\pi_2^{-1}(\G)}}_{ \subseteq \sigma(\mathcal{H})}
			 \subseteq \sigma(\mathcal{H}) \enskip\text{,}
		\end{equation*}
		also auch $\sigma(\mathcal{E}) \subseteq \sigma(\mathcal{H})$.
	\end{itemize}

	Bleibt Nachweis von ($\star\star$):
	
	\begin{itemize}
		\item[$\subseteq$:] ist klar mit $\mathcal{O}$ und $\mathcal{O}'$ einelementig (gilt auch ohne Separabilität)
		\item[$\supseteq$:] Gemäß \cref{satz: 2.9} existiert abzählbare Basis $\G_0$ von $\G$. Seien $\mathcal{O}, \mathcal{O}* \subseteq \G$. Sei
		\begin{align*}
			G^\ast &=\bigcup_{\begin{subarray}{c}
					G \in \mathcal{O}\\
					G' \in \mathcal{O}'
			\end{subarray}}G \times  G'
			=
			\bigcup_{\begin{subarray}{c}
					G,G'\text{ offen}\\
					G \times G' \subseteq G^\ast
			\end{subarray}}
			G \times  G' 
			\overset{(!)}{=}
			\bigcup_{\begin{subarray}{c}
					G_0,G_0' \in \G_0\\
					G \times  G_0' \subseteq  G^\ast
			\end{subarray}}
			G_0 \times  G_0'
		\end{align*}
		eine abzählbare Vereinigung, da $\G_0$ Basis ist, also abzählbar. Somit gilt dann $G^\ast \in \sigma (\{G \times  G' : G,G' \in \G \})$.
	\end{itemize}
\end{proof}

\begin{definition} \label{definition: 3.4}
	Sei $(\Omega,\A)$ ein Messraum.
	Eine Abbildung $X \colon \Omega \to \S$, die $\A$-$\B(\S)$-messbar ist, heißt \begriff{Zufallsvariable} (ZV) in dem metrischen Raum $(\S,d)$ über $(\Omega,\A)$.

	Sei $\P$ ein Wahrscheinlichkeitsmaß auf $(\Omega,\A)$, also $(\Omega,\A,\P)$ ein Wahrscheinlichkeitsraum.
	Das Bildmaß
	\begin{align*}
		\P \circ  X^{-1} &\defeq \P_X \defeq \L(X) \defeq \L(X\,|\;\P) \\
		\brackets{\P \circ X^{-1}}(B)
		&\defeq \P\brackets{X^{-1}(B)}
		=\P\brackets{\menge{\omega \in \Omega : X(\omega) \in  B}}
		\defqe \P(X \in  B) \qquad\forall B \in \B(\S)
	\end{align*}
	heißt \begriff{Verteilung} von $X$ unter $\P$.
\end{definition}

\begin{satz} \label{satz: 3.5}
	Sei $(\S,d)$ ein separabler metrischer Raum und seien $X,Y$ Zufallsvariablen in $(\S,d)$ über $(\Omega,\A)$.
	Dann ist $d(X,Y)$ eine reelle Zufallsvariable.
\end{satz}

\begin{proof}[gute Prüfungsfrage]
	Die Abbildungen $X,Y \colon (\Omega,\A) \to (\S,\B(\S))$ sind messbar genau dann, wenn die Abbildung $(X,Y) \colon (\Omega,\A) \to \big( \S \times \S,\underbrace{\B(\S) \otimes \B(\S)}_{= \B(\S \times \S)}\big)$ messbar ist.
	Jede Metrik ist bekanntlich stetig, also auch $d \colon \brackets{\S \times \S , \G(\S \times \S)} \to \R$\footnote{$\R$ sei mit der natürlichen Topologie versehen}.
	Dann folgt aus \cref{lemma: 3.2}, dass $d \colon \B(\S \times \S) \to \B(\R)$	messbar ist. Damit folgt die Behauptung, denn $d(X,Y) = d \circ (X,Y)$ ist messbar als Komposition von messbaren Abbildungen.
\end{proof}


	% This work is licensed under the Creative Commons
% Attribution-NonCommercial-ShareAlike 4.0 International License. To view a copy
% of this license, visit http://creativecommons.org/licenses/by-nc-sa/4.0/ or
% send a letter to Creative Commons, PO Box 1866, Mountain View, CA 94042, USA.

\section{Fast sichere Konvergenz}
\begin{definition} \label{definition: 3.6}
	Seien $X, X_n$ ($n \in \N$) Zufallsvariablen in einem separablen, metrischen Raum $(\S,d)$ über $(\Omega,\A,\P)$.
	\begin{equation*}
		X_n \ntoinf X \quad \P\text{-fast sicher}
		\defequiv
		\P\bigg( \underbrace{\menge{\omega \in \Omega : d\brackets{X_n(\omega),X(\omega)} \ntoinf 0}}_{\defqe M} \bigg) = 1
	\end{equation*}
\end{definition}

Beachte: Die Definition von $M$ mengentheoretisch aufgeschrieben (Schnitt $\leadsto$ \enquote{für alle}; Vereinigung $\leadsto$ \enquote{Es gibt}):
\begin{equation*}
	M =
	\bigcap_{0 < \epsilon \in \Q}
	\bigcup_{m \in \N}
	\bigcap_{n \ge m}
	\Big\{ \underbrace{d(X_n,X)}_{\defqe \zeta_n} < \epsilon \Big\} \overset{\text{\cref{satz: 3.5}}}{\in} \A
	\text{, denn }\zeta_n^{-1} \brackets{(-\infty,\epsilon)} \in  \A
\end{equation*}
Hierbei gilt $M \in \A$ aufgrund der Abzählbarkeit der beteiligten Schnitte und Vereinigungen. 
Man schreibt auch:
\begin{align*}
	\P\brackets{ \lim_{n \to \infty} X_n = X} = 1 \quad \text{ oder } \quad d(X_n, X) \to  0 \enskip (n \to \infty) \enskip \text{ $\P$-f.\,s.}
\end{align*}

Die bekannten Regeln (Ergebnisse) für \textit{reelle} Zufallsvariablen lassen sich mühelos verallgemeinern; so zum Beispiel der folgende Satz.
\begin{satz} \label{satz: 3.7}
	Der fast-sichere Grenzwert ist $\P$-fast sicher eindeutig:
	\begin{equation*}
		X_n \ntoinf X \quad \P\text{-fast sicher } \und 
		X_n \ntoinf X' \quad\P\text{-fast sicher }
		\follows 
		X = X' \quad \P\text{-fast sicher}
	\end{equation*}
\end{satz}
\begin{proof} 
	Es ist $\menge{ X\neq X'} \subseteq \menge{X_n \not\to  X} \cup \menge{X_n \not\to  X'}$ und $\P(X_n \not\to X) + \P(X_n \not\to X') = 0+0$. Also ist auch $\P(X \neq X') = 0$ und mit dem Gegenereignis gilt $\P\brackets{X = X'} = 1-\P\brackets{X\neq X'} = 1$.
\end{proof}

\begin{satz} \label{satz: 3.8}
	Seien $X, X_n$ ($n \in \N$) Zufallsvariablen im separablen, metrischen Raum $(\S,d)$ und sei $f \colon (\S,d) \to (\S',d')$ $\B_d(S)$ - $\B_{d'}(S')$-messbar und stetig in $X$ $\P$-fast sicher.
	\footnote{\enquote{stetig in $X$ $\P$-fast sicher} bedeutet,
		dass $\P(\menge{\omega  \in  \Omega : f \text{ in } X(\omega) \text{ stetig}}) = 1$. Insbesondere ist $\menge{\omega  \in  \Omega : f \text{ in } X(\omega) \text{ stetig}}$ messbar, was keine Selbstverständlichkeit ist.}
	Dann gilt:
	\begin{equation*}
		X_n \ntoinf	X \quad \P\text{-fast sicher }
		\follows f(X_n) \ntoinf f(X) \quad \P\text{-fast sicher}
	\end{equation*}
\end{satz}

\begin{proof}
	Aufgrund der Folgenstetigkeit gilt
	\begin{equation*}
		\menge{X_n \ntoinf  X} \cap \menge{f \text{ stetig in }X}
		\subseteq
		\menge{f(X_n) \ntoinf f(X)}
	\end{equation*}
	Da abzählbare Schnitte von Einsmengen stets Einsmengen sind, d.\,h.
	\begin{equation*}
		\forall i \in \N : \, \P(E_i)=1
		\follows \P\brackets{\bigcap_{i \in \N} E_i}=1
		\qquad \forall \menge{E_i} \subseteq \Omega \mit \P(E_i)=1
	\end{equation*}
	folgt $1 = \P\brackets{\menge{X_n \ntoinf X \text{ und $f$ stetig in } X}} \le\P\brackets{\menge{f(X_n) \ntoinf  f(X)}} \le 1$ und somit schon $\P\brackets{\menge{f(X_n) \ntoinf  f(X)}} = 1$.
\end{proof}

%\begin{satz}[Konvergenz-Kriterium]\label{Satz3.9}
%	\begin{align*}
%		X_n
%		\ntoinf
%		% CHECKED: '\longrightarrow' used.
%		X\quad\P\text{-fast sicher}
%		\Leftrightarrow
%		\forall\epsilon>0:\limn\P\brackets{\sup_{m\ge n} d(X_m,X)>\epsilon}=0
%	\end{align*}
%\end{satz}
%\begin{proof}
%	Man ersetze im Beweis für den Fall reeller Zufallsvariablen $\abs{X_n-X}$ durch $d(X_n,X)$.
%	Und beachte, dass alle Schlussfolgerungen bestehen bleiben. Hier ausführlich:
%	\begin{align*}
%		X_n \overset{n\to \infty}{\longrightarrow} X~\P \text{-fast sicher}
%		&\Leftrightarrow 1 = \P\brackets{\bigcup_{0 < \epsilon} \bigcap_{m  \in  \N} \bigcup_{n \ge m} d(X_n, X) < \epsilon} \\
%		&\Leftrightarrow 0 = \P\brackets{\bigcap_{0 < \epsilon} \bigcup_{m  \in  \N} \bigcap_{n \ge m} d(X_n, X) \ge \epsilon} \\
%		&\Leftrightarrow \forall \epsilon > 0: 0 = \P\brackets{\bigcup_{m  \in  \N} \bigcap_{n \ge m} d(X_n, X) \ge \epsilon}
%	\end{align*}
%	Da für $m \to  \infty$ die Mengen $\menge{\bigcap_{n \ge m} d(X_n, X) \ge \epsilon}$ kleiner werden,
%	ist dies äquivalent zu
%	\begin{align*}
%		\forall \epsilon > 0\, \forall δ > 0 \,∃m  \in  \N: \P\brackets{\bigcap_{n \ge m} d(X_n, X) \ge \epsilon} &< δ \\
%		\Leftrightarrow \forall \epsilon > 0: \lim_{m \to  \infty} \P\brackets{\bigcap_{n \ge m} d(X_n, X) \ge \epsilon} &= 0 \\
%		\Leftrightarrow \forall \epsilon > 0: \lim_{m \to  \infty} \P\brackets{\sup_{n \ge m} d(X_n, X) \ge \epsilon} &= 0
%		\qedhere
%	\end{align*}
%\end{proof}
%Ein sehr nützliches Kriterium ist Folgendes:
%\begin{satz}\label{Satz3.10}
%	\begin{align*}
%		\sum_{n \in \N_{\ge1}}\P\brackets{d(X_n,X)>\epsilon} < \infty \qquad \forall \epsilon>0
%		\Rightarrow  X_n
%		\ntoinf
%		% CHECKED: '\longrightarrow' used.
%		X\quad\P\text{-fast sicher}
%	\end{align*}
%	\begin{bemerkung}
%		% noNumber
%		Die $\P(d(X_n, X) > \epsilon)$ werden in der Statistik \define{Fehlerwahrscheinlichkeit}
%		oder \define{tail probability} genannt.
%		Um Fehlerwahrscheinlichkeiten abzuschätzen, also die Voraussetzung für diesen Satz für einen speziellen Fall zu zeigen, nutzt man häufig sogenannte Maximalungleichungen wie die Markov-Ungleichung und die Tschebychew-Ungleichung.
%	\end{bemerkung}
%\end{satz}
%
%\begin{proof}
%	Setze
%	\begin{align*}
%		A_n(\epsilon):=\menge{d(X_n,X)>\epsilon}\overset{\ref{Satz3.5}}{ \in }\A
%	\end{align*}
%	Dann folgt aus dem \undefine{ersten Borel-Cantelli-Lemma}:
%	\begin{align*}
%		&\P\brackets{\limsup_{n\to\infty} A_n(\epsilon)}=0\qquad\forall\epsilon>0
%	\end{align*}
%		Mit
%	\begin{align*}
%		\liminf_{n\to\infty}\brackets{A_n(\epsilon)^C}
%		\overset{\text{Def}}=
%		\bigcup_{m \in \N}\bigcap_{n\ge m}\brackets{A_n(\epsilon)}^C
%		&= \brackets{\bigcap_{m \in \N}\bigcup_{n\ge m}\menge{ d(X_n,X)>\epsilon}}^C \\
%		&\overset{\Def}= \brackets{\limsup_{n \to  \infty} A_n(\epsilon)}^C
%	\end{align*}
%	folgt dann
%	\begin{align*}
%		1=\P\brackets{\brackets{\limsup_{n\to\infty} A_n(\epsilon)}^C}
%		=\P\brackets{\liminf_{n\to\infty}\brackets{A_n(\epsilon)^C}}\qquad\forall\epsilon>0
%	\end{align*}
%	Da abzählbare Durchschnitte von Eins-Mengen (also Mengen mit $\P$-Maß 1) wieder Eins-Mengen sind, folgt schließlich:
%	\begin{align*}
%		\P\brackets[\Bigg]{\underbrace{
%			\bigcap_{0<\epsilon \in \Q}\bigcup_{m \in \N}\bigcap_{n\ge m}\menge{d(X_n,X)\le\epsilon}
%		}_{\menge{ X_n\to X}=\menge{ d(X_n,X)\to0}}}=1
%		&
%		\qedhere
%	\end{align*}
%\end{proof}
%
%Weitere Eigenschaften der fast sicheren Konvergenz von Zufallsvariablen in metrischen Räumen finden sich z.\,B.\ in \cite[Kapitel 8.2]{gaensslerstute1977Wahrscheinlichkeitstheorie}.
%

	% This work is licensed under the Creative Commons
% Attribution-NonCommercial-ShareAlike 4.0 International License. To view a copy
% of this license, visit http://creativecommons.org/licenses/by-nc-sa/4.0/ or
% send a letter to Creative Commons, PO Box 1866, Mountain View, CA 94042, USA.

\section{Stochastische Konvergenz}

\begin{definition} \label{definition: 3.11}
	$X_n$ \begriff{konvergiert stochastisch} bzw. \begriff{in Wahrscheinlichkeit} gegen $X$, in Symbolen:
	\begin{equation*}
		X_n \probto X \defequiv \forall \epsilon > 0:
		\P\brackets{\menge{d(X_n,X) > \epsilon}} \ntoinf 0
	\end{equation*}
	d.\,h.\ $d(X_n, X) \probto 0$.
\end{definition}

\begin{satz} \label{satz: 3.12}
	$\displaystyle X_n \ntoinf X \quad \P\text{-fast sicher} \follows X_n \probto X$
\end{satz}

\begin{proof}
	Gemäß \cref{satz: 3.9} gilt
	\begin{equation*}
		\forall \epsilon > 0 : \enskip 
		0 \le \P\brackets{d(X_n,X) > \epsilon}
		\le \P\brackets{\sup_{m \ge n} d(X_m,X) > \epsilon}
		\ntoinf	0
	\end{equation*}
\end{proof}

Die Umkehrung von \cref{satz: 3.12} gilt im Allgemeinen \textit{nicht}, aber es gilt das folgende Teilfolgenkriterium.

\begin{satz}[Teilfolgenkriterium für stochastische Konvergenz] 
	\label{satz: 3.13}
	Folgende Aussagen sind äquivalent:
	\begin{enumerate}[label=(\arabic*)]
		\item \label{it: 3.13WahrKonv} $\begin{aligned} X_n \probto	X \end{aligned}$
		\item \label{it: 3.13TFfsKonv} Zu jeder Teilfolge (TF) $(X_{n'})$ von $(X_n)_{n \in \N}$ existiert eine Teilfolge $(X_{n''})$ von $(X_{n'})$ derart, dass $X_{n''} \overset{n''\to \infty}{\longrightarrow} X$ $\P$-fast sicher.
	\end{enumerate}
\end{satz}

Die Notation $(X_{n''})$ stammt aus \cite{billingsley2008probability}.
%Der folgende Beweis nutzt sie nicht, da er nicht aus der Vorlesung, sondern
%aus der Feder des Mitschreibers stammt, der sie unsauber findet.

\begin{proof}
	Der Beweis verläuft wie in $\R$.
	\begin{description}
		\hinrichtung Sei $\epsilon > 0$ beliebig und $(X_{n_k})_k$ eine Teilfolge von $(X_n)_n$.
		Da $\P\brackets{\menge{d(X_n, X) > \epsilon}} \to 0$, können wir eine Teilteilfolge $(X_{n_{k_i}})_i \defeq (Y_i)_i$ finden mit $\P\brackets{\menge{d(Y_i, X) > \epsilon}} < i^{-2}$. Dann ist für alle $n \in \N$ und mit $n \to \infty$
		\begin{align*}
			\P\brackets{\sup_{i \ge n} d(Y_i, X) < \epsilon}
			= \P\brackets{\bigcup_{i \ge n} \menge{d(Y_i, X) < \epsilon}}
			\leq \sum_{i = n}^{\infty} \P\brackets{\menge{d(Y_i, X) < \epsilon}}
			\leq \sum_{i = n}^{\infty} i^{-2} \to 0 
		\end{align*}
		Laut dem Konvergenz-Kriterium in \cref{satz: 3.9} bedeutet das $Y_i \to X$ $\P$-fast sicher.
		
		\rueckrichtung Wir nehmen an, es gäbe ein $\epsilon > 0$ und eine Teilfolge $(X_{n_k})_k$ der Folge $(X_n)_n$, sodass $\P\brackets{d(X_{n_k}, X) > \epsilon} > \delta > 0$ für ein passendes $\delta$.
		Dann konvergiert nach \cref{it: 3.13TFfsKonv} eine Teilteilfolge $(X_{n_{k_i}})_i$ $\P$-fast sicher gegen $X$, also laut \cref{satz: 3.12} auch stochastisch. 
		Das widerspricht aber der obigen Annahme über $(X_{n_k})_k$, die sich
		auch auf $(X_{n_{k_i}})_i$ überträgt.
	\end{description}
	~
\end{proof}

Mit dem Teilfolgenkriterium lassen sich Rechenregeln für fast sichere Konvergenz auf stochastische Konvergenz übertragen.

\begin{korollar} \label{korollar: 3.14}
	\begin{enumerate}[label=(\alph*)]
		\item \label{it: 3.14probConvEq} $X_n \probto X$ und $X_n \probto X'
			\follows X = X'\quad\P\text{-fast sicher}$
		\item \label{it: 3.14probConvCont} Gelte $X_n \probto X$ in $(\S,d)$ und $f \colon (\S,d) \to (\S',d')$ messbar mit $f$ stetig in $X$ $\P$-fast sicher. Dann gilt auch $f(X_n) \probto f(X)$.
	\end{enumerate}
\end{korollar}

\begin{proof}
	\begin{enumerate}[label=(zu \alph*), leftmargin=*]
		\item Wegen $X_n \probto X$ existiert nach \cref{satz: 3.13} eine Teilfolge $(X_{n'}) \subseteq(X_n)_{n \in \N}$ mit $X_{n'}
		\overset{n' \to \infty}{\longrightarrow} X$ $\P$-fast sicher.
		Zu $(X_{n'})$ existiert (wegen $X_n\probto X'$ und \cref{satz: 3.13}) eine Teilfolge $(X_{n''})$ von $(X_{n'})$, sodass aus $X_{n''}	\ntoinf  X' \enskip \P\text{-fast sicher}$ mit \cref{satz: 3.7} schon $X = X'$ fast sicher folgt.
		\item Zu einer beliebigen Teilfolge $(X_{n'})$ von $(X_n)$ existiert eine Teilfolge $(X_{n''})$ von $(X_{n'})$ mit $X_{n''} \to X$ $\P$-fast sicher.
		Mit \cref{satz: 3.8} folgt $f(X_{n''}) \to f(X)$ $\P$-fast sicher und mit \cref{satz: 3.13} angewendet auf $(f(X_n))_n$ gilt $f(X_n) \probto f(X)$.
	\end{enumerate}
\end{proof}
	% This work is licensed under the Creative Commons
% Attribution-NonCommercial-ShareAlike 4.0 International License. To view a copy
% of this license, visit http://creativecommons.org/licenses/by-nc-sa/4.0/ or
% send a letter to Creative Commons, PO Box 1866, Mountain View, CA 94042, USA.

\section{Konvergenz in Produkträumen}

Seien $(\S,d)$ und $(\S',d')$ separable metrische Räume. Dann ist auch $(\S \times \S',d \times  d')$ ein separabler metrischer Raum.
Dies folgt z.\,B.\ aus dem \textit{Satz von der koordinatenweise Konvergenz}:

\begin{equation}
	\label{eq: 3.1KoordinatenweissKonvergenz}
%	\tag{3.1}
	\brackets{a_n, a_n'} \overset{d \times d'}{\longrightarrow} (a, a')
	\equivalent
	(a_n) \overset{d}{\longrightarrow} a
	\und
	(a_n') \overset{d'}{\longrightarrow} a'
\end{equation}

Es gibt auch eine \enquote{stochastische Versionen} dieses Satzes.

\begin{satz}[Koordinatenweise Konvergenz vom Zufallsgrößen]
	\label{satz: 3.15} ~
	\begin{enumerate}[label=(\arabic*)]
		\item \label{it: 3.15koordKonvfs} 
			$(X_n,X_n')
			\overset{d \times d'}{\longrightarrow}
			(X,X') \enskip ~\P\text{-f.\,s.}
			\equivalent
			X_n \overset{d}{\longrightarrow} X \enskip \P\text{-f.\,s.\ } 
			\enskip \und \enskip
			X_n' \overset{d'}{\longrightarrow} X' \enskip \P\text{-f.\,s.}$
		\item \label{it: 3.15koordKonvprob} $(X_n,X_n') \probto (X,X')
			\equivalent
			X_n \probto X \enskip \und \enskip X_n' \probto X'$
	\end{enumerate}
\end{satz}

\begin{proof}
	\begin{enumerate}[label=(zu \arabic*), leftmargin=*]
		\item Nach \cref{eq: 3.1KoordinatenweissKonvergenz} folgt aus $(X_n,X_n') \overset{d \times d'}{\longrightarrow} (X,X') \enskip \P$-fast sicher auch $X_n \to X$ und $X_n' \to X' \enskip \P$-fast sicher. Da der Schnitt von Eins-Mengen wieder eine Eins-Menge ist, ist dies wiederum äquivalent zu $X_n \to X \enskip \P$-fast sicher und $X_n' \to X' \enskip \P$-fast sicher. \textit{Achtung} -- die \enquote{Stellung} der Fast-Sicherheit spielt eine Rolle!
		\item Die linke Seite ist wegen \eqref{eq: 3.1KoordinatenweissKonvergenz} und \cref{satz: 3.13} äquivalent dazu, dass für alle Teilfolgen $(X_{n'},X_{n'}') \subseteq (X_n,X_n')$ eine Teilteilfolge $(X_{n''},X_{n''}')$ existiert mit $(X_{n''},X_{n''}') \ntoinf 0$ fast sicher.
		Also gilt wegen \cref{it: 3.15koordKonvfs}
		\begin{align*}
			X_{n''}\to X\text{ f.\,s.\ und }X_{n''}'\to X'\text{ f.\,s.}
		\end{align*}
		Somit existiert für alle Teilfolgen $(X_{n'}) \subseteq (X_n)$ eine Teilteilfolge $(X_{n''}) \subseteq (X_{n'})$ mit $X_{n''} \to X$ fast sicher. Dies ist wegen \eqref{eq: 3.1KoordinatenweissKonvergenz} äquivalent zu $X_n \probto X$. Analog gilt auch $X_n'\to X'$.
	\end{enumerate}
\end{proof}
	% This work is licensed under the Creative Commons
% Attribution-NonCommercial-ShareAlike 4.0 International License. To view a copy
% of this license, visit http://creativecommons.org/licenses/by-nc-sa/4.0/ or
% send a letter to Creative Commons, PO Box 1866, Mountain View, CA 94042, USA.

\subsection{Gleichheit in Verteilung}
\begin{definition} \label{definition: 3.16}
	Seien $X_i$ ($i = 1,2$) Zufallsvariablen in $(\S,d)$ über $(\Omega_i,\A_i,\P_i)$. Sie heißen \begriff{gleich in Verteilung}, in Zeichen
	\begin{align*}
		X_1 \disteq X_2 \defequiv \P_1 \circ X_1^{-1} = \P_2 \circ X_2^{-1}
	\end{align*}
\end{definition}

Wir wollen Verteilungsgleichheit in folgendem Satz charakterisieren.

\begin{satz} \label{satz: 3.17}
	Es gilt:
	\begin{enumerate}[label=(\alph*)]
		\item \label{it: 3.17f} Seien $P$ und $Q$ Wahrscheinlichkeitsmaße auf $\B(\S)$. Dann gilt:
		\begin{equation*}
			P \equiv Q
			\equivalent
			\int f \diffskip P = \int f \diffskip Q \qquad \forall f \in C^b(\S) \text{ glm.\ stetig}
		\end{equation*}
		\item \label{it: 3.17E} Es gilt
		\begin{equation*}
			X \disteq Y \equivalent
			\EW[f(X)] = \int_{\Omega_1} f(X_1)  \diffskip{\P_1} = \int_{\Omega_2} f(X_2)  \diffskip{\P_2} = \EW[f(Y)]
		\end{equation*}
		für alle $f \in C^b(\S)$ gleichmäßig stetig.
	\end{enumerate}
\end{satz}

\begin{proof}
	\begin{enumerate}[label=(zu \alph*), leftmargin=*]
		\item \begin{description}
			\hinrichtung Klar.
			\rueckrichtung Aus \cref{lemma: 3.2}, \cref{it: 3.2-OpenClosed} wissen wir, dass $\B(\S) = \sigma\brackets{\mathcal{F}(\S)}$ und $\F$ ist durchschnittsstabil. Wegen des \textit{Maßeindeutigkeitssatz} reicht es zu zeigen, dass $\P(F) = \Q(F)$ für alle $F \in \mathcal{F}(S)$. Sei nun $F \subseteq \S$ abgeschlossen. Setze $f_k (x) \defeq \phi \brackets{k * d(x,F)}$ (vgl.\ \cref{satz: 2.4}, $\phi$ wie dort). Aus \cref{lemma: 2.3} folgt, dass die $f_k$ beschränkt und gleichmäßig stetig sind mit $f_k \searrow \one_F$ für $k \to \infty$. Also gilt:
			\begin{align*}
				\P(F) 
				&= \int \one_F \diffskip{\P}
				= \int \lim_{k\to\infty} f_k  \diffskip\P
				\intertext{und mit monotoner Konvergenz}
				&= \lim_{k \to \infty} \int f_k \diffskip\P
				= \lim_{k \to \infty} \int f_k \diffskip\Q
				\intertext{Wiederum mit monotoner Konvergenz}
				&= \int \lim_{k \to \infty} f_k \diffskip\Q
				= \int \one_F \diffskip\Q
				= \Q(F)
			\end{align*}
			Da $F$ beliebig war, folgt die Behauptung.
		\end{description}
		\item Dies folgt aus \cref{it: 3.17f} mit dem Transformationssatz \eqref{eqTrafo}:
			\begin{align*}
			X \disteq Y
			&\overset{\ref{definition: 3.16}}{\Leftrightarrow}
			\P \circ X^{-1} = \P \circ Y^{-1} \\
			&\overset{\ref{it: 3.17f}}{\Leftrightarrow}
			\int_{\S} f \diffskip{(\P\circ X^{-1})} = \int_{\S} f \diffskip{(\P\circ Y^{-1})}
			&\forall f \in C^b(\S) \text{ glm. stetig} \\
			&\Leftrightarrow
			\EW[f(X)] = \EW[f(Y)]
			&\forall f \in C^b(\S) \text{ glm. stetig}
		\end{align*}
		Verwende dafür im letzten Schritt den Transformationssatz, d.h.
		\begin{equation*}
			\int_{\S} f \diffskip{(\P\circ X^{-1})}
			\overset{\text{Trafo}}{=}
			\int_{\Omega} f \circ X \diffskip\P
			\overset{\text{Def}}{=}
			\EW[f \circ X] = \EW[f(X)]
		\end{equation*}
	\end{enumerate}
\end{proof}
	
	\chapter{Verteilungskonvergenz von Zufallsvariablen in metrischen Räumen}
	% This work is licensed under the Creative Commons
% Attribution-NonCommercial-ShareAlike 4.0 International License. To view a copy
% of this license, visit http://creativecommons.org/licenses/by-nc-sa/4.0/ or
% send a letter to Creative Commons, PO Box 1866, Mountain View, CA 94042, USA.

\section{Satz von Portmanteau}

Seien $X$ und $(X_n)_{n \in \N}$ Zufallsvariablen in $(\S,d)$ über $(\Omega,\A,\P)$.
Dann sind
\begin{equation*}
	P \defeq \P \circ X^{-1} \qquad \und \qquad P_n \defeq \P \circ X_n^{-1} \qquad (n \in \N)
\end{equation*}
Wahrscheinlichkeitsmaße auf $\B(\S)$.

\begin{definition} \label{definition: 4.1}
	\begin{enumerate}[label=(\arabic*)]
		\item \label{it: 4.1Pweak} Seien $P$ und $(P_n)_{n \in \N}$ Wahrscheinlichkeitsmaße auf $\B(S)$.
		Dann \begriff{konvergiert $P_n$ schwach} gegen $P$, in Zeichen
		\begin{equation*}
			P_n \weakto P \defequiv \int f \diffskip{P_n} \ntoinf \int f \diffskip P \qquad \forall f \in C^b(\S)
		\end{equation*}
		\item \label{it: 4.1Xweak} $X_n$ \begriff{konvergiert in Verteilung} gegen $X$ in Raum $(\S,d)$, in Zeichen
			\begin{align*}
				X_n \distto X \text{ in } (\S,d) \defequiv \P \circ X_n^{-1} \weakto \P \circ X^{-1}
			\end{align*}
			Das $\text{d}$ steht für \enquote{distribution}. Alternative Schreibweise: $X_n \overset{\mathcal{L}}{\longrightarrow} X$.
	\end{enumerate}
\end{definition}

Der folgende Satz gibt eine äquivalente Charakterisierung von $\weakto$ bzw.\ $\distto$.

% This work is licensed under the Creative Commons
% Attribution-NonCommercial-ShareAlike 4.0 International License. To view a copy
% of this license, visit http://creativecommons.org/licenses/by-nc-sa/4.0/ or
% send a letter to Creative Commons, PO Box 1866, Mountain View, CA 94042, USA.

\begin{theorem}[Portmanteau-Theorem]\label{satz: 4.2}
	Folgende Aussagen sind äquivalent:
	\begin{enumerate}[label=(\arabic*)]
		\item \label{it:4.2weakConv} $\begin{aligned}
			P_n \weakto P
		\end{aligned}$
	\item \label{it:4.2f} $\begin{aligned}
			\int f \diffskip P_n \longrightarrow \int f \diffskip P \qquad \forall f \in C^b(\S) \text{ glm. stetig}
			% Standardfrage für die Prüfung: was sagt das Portmanteau-Theorem und was ist schwache Konvergenz?
		\end{aligned}$
	\item \label{it:4.2ClosedSets} $\begin{aligned}
			\limsup_{n \to \infty} P_n(F) \le P(F) \qquad \forall F \in \F(\S)
		\end{aligned}$
	\item \label{it:4.2OpenSets} $\begin{aligned}
			\liminf_{n \to \infty} P_n(G)\ge P(G) \qquad \forall G \in \G(\S)
		\end{aligned}$
	\item \label{it:4.2BorelSets} $\lim\limits_{n \to \infty} P_n(B) = P(B) \qquad \forall B \in \B(\S) \mit P(\underbrace{\partial B}_{ \in\F(\S)}) = 0$ \\
		Die Mengen $B \in \B(\S)$ mit $P(\partial B)=0$ heißen \begriff{$P$-randlos}.
	\end{enumerate}
\end{theorem}

\begin{proof}
	\begin{description}
		\item[\ref{it:4.2weakConv} $\Rightarrow$ \ref{it:4.2f}:] 
		Folgt aus \cref{definition: 4.1}.
		%
		\item[\ref{it:4.2f} $\Rightarrow$ \ref{it:4.2ClosedSets}:] 
		Sei $F \in\F(\S)$. Der Beweis von \cref{satz: 3.17} zeigt, dass es eine Folge $(f_k)_{k \in \N}$ von gleichmäßig stetigen und beschränkten Funktionen auf $\S$ gibt mit $f_k \downarrow \one_F$. Dann gilt:
		\begin{equation*}
			\limsup_{n \to \infty} P_n(F)
			= \limsup_{n \to \infty} \int \underbrace{\one_F}_{\le f_k} \diffskip P_n
			\overset{\text{Mon.}}{\le}
			\limsup_{n \to \infty} \int f_k \diffskip P_n
			\overset{\text{Vor.}}{=}
			\int f_k \diffskip P \quad \forall k \in \N 
			\satzende
		\end{equation*}
		Mit monotoner Konvergenz folgt daraus $\int f_k \diffskip P \longrightarrow \int \one_F \diffskip  P = P(F)$ und somit für $k \to \infty$ auch \ref{it:4.2ClosedSets}.
		%
		\item[\ref{it:4.2ClosedSets} $\Leftrightarrow$ \ref{it:4.2OpenSets}:]
		Nutze den Übergang zum Komplement sowie die Rechenregeln für $\liminf$ und $\limsup$. Sei $G$ offen, das heißt $G \in \G(\S)$ und damit $G^\complement \in \F(\S)$. Dann gilt
		\begin{equation*}
			\liminf_{n \to \infty} P_n(G)
			= \liminf_{n \to \infty} \brackets{1-\P_n(G^\complement)}
			= 1 - \underbrace{\limsup_{n \to \infty} P_n(G^\complement)}_{\le P(G^\complement)}
			\ge 1 - P(G^\complement)
			= P(G) 
			\satzende
		\end{equation*}
	\end{description}

%	\paragraph{Zeige \ref{it:4.2ClosedSets} $ \Rightarrow $ \ref{it:4.2weakConv}:}
%	Sei $f \in C^b(\S)$ beliebig. Zeige zunächst:
%	\begin{align}\label{eqProof1.4.2Sternchen}\tag{$\ast$}
%		\limsup_{n \to \infty}\int f\diffskip \P_n\le\int f\diffskip \P
%	\end{align}
%	\subparagraph{1. Schritt:} Sei $0\le f<1$. Setze
%	\begin{align*}
%		F_i:=\set{ f\ge\frac{i}{k}}=\set{ x \in\S:f(x)\ge\frac{i}{k}}
%		= f^{-1}\klammern{\intervallHO{\frac1k}{∞}},\qquad  \forall\; 0\le i\le k,\;k \in\N
%	\end{align*}
%	($F_i$ hängt auch von $k$ ab, aber diese Abhängigkeit wird in der Notation nicht
%	demonstriert.)
%	Dann gilt $F_i \in\F~ \forall i$, da $f$ stetig. Da
%	\begin{align*}
%		\int_{\S}f\diffskip \P
%		\stackeq{\text{Lin}}
%		\sum_{i=1}^k\int\indi_{\set{\frac{i-1}{k}\le f<\frac{i}{k}}}· f\diffskip \P
%	\end{align*}
%	folgt wegen Monotonie des Integrals
%	\begin{align}\label{eqProof1.4.2Plus}\tag{+}
%		\sum_{i=1}^k\underbrace{\frac{i-1}{k}}_{=\frac{i}{k}-\frac{1}{k}}·
%	\P\klammern{\frac{i-1}{k}\le f<\frac{i}{k}}
%		\le
%		\int f\diffskip \P
%		\le
%		\sum_{i=1}^k \frac{1}{k}· \P\klammern[\Bigg]{\underbrace{\frac{i-1}{k}\le f<\frac{i}{k}}_{F_{i-1}\setminus F_i}}
%	\end{align}
%	Die rechte Summe in \eqref{eqProof1.4.2Plus} ist gleich
%	\begin{align*}
%		&\frac{1}{k}·\sum_{i=1}^k i·\klammern[\big]{ \P(F_{i-1})-\P(F_i)}
%		\quad \text{da $F_i ⊂ F_{i-1}$}\\
%		&=\frac{1}{k}·\Big(\P(F_0)-\P(F_1)+2· \P(F_1)-2· \P(F_2)+3· \P(F_2)-3· \P(F_3)+\\
%		% CHECKED: '\Big' used.
%		&\qquad+…+(k-1)· \P(F_{k-2})-(k-1)· \P(F_{k-1})+k· \P(F_{k-1})-k· \P(F_k)\Big)\\
%		% CHECKED: '\Big' used.
%		&=\frac{1}{k}·\klammern[\Big]{\underbrace{\P(F_0)}_{=1}+\P(F_1)+\P(F_2)+…+\P(F_{k-1})- k· \underbrace{\P(F_k)}_{=0}} \quad (0 \le f < 1)\\
%		&=\frac{1}{k}+\frac{1}{k}·\sum_{i=1}^{k-1} \P(F_i)
%	\end{align*}
%	Da die linke Summe in \eqref{eqProof1.4.2Plus} gleich der rechten Summe in \eqref{eqProof1.4.2Plus} minus $\frac{1}{k}$ ist, folgt
%	\begin{align}\label{eqProof1.4.2DoppelSternchen}\tag{$\ast\ast$}
%		\sum_{i=1}^{k-1} \P(F_i)
%		\le\int f\diffskip P
%		\le\frac{1}{k}+\frac 1k \sum_{i=1}^{k-1} \P(F_i)
%	\end{align}
%	Beachte, \eqref{eqProof1.4.2DoppelSternchen} gilt für \emph{jedes} Wahrscheinlichkeitsmaß $\P$, also auch für $\P_n$.
%	Damit folgt:
%	\begin{align*}
%		\limsup_{n \to \infty}\int f\diffskip \P_n
%		\overset{\eqref{eqProof1.4.2DoppelSternchen}}&{\le}
%		\frac{1}{k}+\frac 1k \sum_{i=1}^{k-1}\underbrace{
%			\limsup_{n \to \infty} \P_n(F_i)
%		}_{\overset{\ref{it:4.2ClosedSets}}{\le}\P(F_i)~ \forall i}\\
%		\overset{\ref{it:4.2ClosedSets}}&{\le}
%		\frac{1}{k}+\underbrace{\frac 1k \sum_{i=1}^{k-1} \P(F_i)}_{
%			\overset{\eqref{eqProof1.4.2DoppelSternchen}}{\le}\int f\diffskip \P
%		}\\
%		\overset{\eqref{eqProof1.4.2DoppelSternchen}}&{\le}
%		\frac{1}{k}+\int f\diffskip \P\qquad \forall k \in\N
%	\end{align*}
%	Grenzwertbildung $k \to \infty$ liefert \eqref{eqProof1.4.2Sternchen}.
%	\paragraph{2. Schritt:} Da $f \in C^b(\S)$ beliebig, gilt wegen Beschränktheit von $f$:
%	\begin{align*}
%		∃ a<b:a\le f<b
%		 \Rightarrow  g(x):=\frac{f(x)-a}{b-a}\text{ ist stetig und } 0\le g<1
%	\end{align*}
%	Daraus folgt
%	\begin{align*}
%		\limsup_{n \to \infty}\int f\diffskip \P_n
%		&=\limsup_{n \to \infty}\int (b-a)· g+a\,\diffskip \P_n\\
%		&=\limsup_{n \to \infty}\klammern{(b-a)·\int g\,\diffskip \P_n+a}\\
%		&\le(b-a)·\underbrace{\limsup_{n \to \infty}\int g\diffskip \P_n}_{\le\int g\diffskip P\text{, wg.\ 1.\ Schritt}}+a\\
%		&\le(b-a)·\int g\diffskip \P + a\\
%		\overset{\text{Lin}}&=
%		\int f\diffskip \P
%	\end{align*}
%	Damit ist \eqref{eqProof1.4.2Sternchen} gezeigt. Übergang zu $-f$ in \eqref{eqProof1.4.2Sternchen} liefert
%	\begin{align*}
%		\liminf_{n \to \infty}\int f\diffskip \P_n
%		&=\liminf_{n \to \infty}-\int -f\diffskip \P_n\\
%		&=-\limsup_{n \to \infty}\int \underbrace{-f}_{ \in C^b(\S)}\diffskip \P\\
%		\overset{\eqref{eqProof1.4.2Sternchen}}&{\ge}
%		-\int -f\diffskip \P\\
%		\overset{\text{Lin}}&=
%		\int f\diffskip \P
%		\qquad \forall f \in C^b(\S)\\
%		 \Rightarrow  \intf \diffskip\P &\le \liminf_{n → ∞} \intf\diffskip\P_n \le \limsup_{n → ∞} \intf\diffskip\P_n \overset{\eqref{eqProof1.4.2Sternchen}}{\le} \intf \diffskip\P \\
%		 \Rightarrow  \lim_{n → ∞} \int f \diffskip\P_n &= \intf \diffskip\P \quad  \forall f  \in C^b(S) \\
%		& \Rightarrow \ref{it:4.2weakConv}
%	\end{align*}
%
%	\subparagraph{Zeige \ref{it:4.2ClosedSets} $ \Rightarrow $ \ref{it:4.2BorelSets}:}
%	Sei $B \in\B(\S)\mit \P(\partial B)=0$. Dann gilt:
%	\begin{align*}
%		\P(\overline{B})
%		\overset{\ref{it:4.2ClosedSets}}&{\ge}
%		\limsup_{n \to \infty} \underbrace{\P_n(\overbrace{\overline{B}}^{⊇ B})}_{\ge \P_n(B)}
%		\ge\limsup_{n \to \infty} \P_n(B)
%		\overset{\text{stets}}{\ge}
%		\liminf_{n \to \infty} \underbrace{\P_n(\overbrace{B}^{⊇\overset{∘}{B}})}_{
%			\ge \P_n(\inner B)
%		}\\
%		\overset{\ref{it:4.2ClosedSets}⇔ \ref{it:4.2OpenSets}}&{\ge}
%		\P\klammern[\big]{\inner B}
%		=\P(\overline{B})
%		=\P(B),
%	\end{align*}
%	denn:
%	\begin{align*}
%		0
%		=\P(\overbrace{\partial B}^{\overline{B}\setminus \inner B})
%		=\P(\overline{B})-\P(\inner B)
%		 \Rightarrow 
%		\P(\inner B)\le \P(B)\le \P(\overline{B})=\P(\inner B)
%	\end{align*}
%	Damit folgt $\liminf=\limsup$ und folglich $\lim_{n \to \infty} \P_n(B)=\P(B)$.
%	\paragraph{Zeige \ref{it:4.2BorelSets} $ \Rightarrow $ \ref{it:4.2ClosedSets}:}
%	Sei $F \in\F$ (abgeschlossen) beliebig. Dann gilt $ \forallε>0:$
%	\begin{align}\label{eqProof1.4.2SternchenUnten}\tag{$\ast$}
%		\partial\,\set{ x \in\S:d(x,F)\leε}
%		⊆\set{ x \in\S:d(x,F)=ε}
%	\end{align}
%	denn: Sei $x \in\partial\set{ x \in\S:d(x,F)\leε}$. Dann gilt:
%	\begin{align*}
%		&∃ (x_n)_{n \in\N}: \forall n \in\N:d(x_n,F)\leε∧ \lim_{n \to \infty} x_n=x\\
%		&∃ (ζ_n)_{n \in\N}: \forall n \in\N:d(ζ_n,F)>ε∧ \lim_{n \to \infty} ζ_n=x
%	\end{align*}
%	Da $d(·,F)$ stetig ist gemäß \ref{lemma2.3} \ref{it:distStetig}, folgt
%	\begin{align*}
%		ε \le \lim_{n → ∞} d(ζ_n, F) = d(x,F) = \lim_{n → ∞} d(x_n, F) \le ε.
%	\end{align*}
%	Wegen \eqref{eqProof1.4.2SternchenUnten} sind
%	\begin{align*}
%		A_ε:=\partial\set{ x \in\S:d(x,F)\leε}\qquad \forallε>0
%	\end{align*}
%	paarweise disjunkt, da bereits die Obermengen paarweise disjunkt sind. Dann folgt
%	\begin{align}\label{eqProof1.4.2DoppelSternchenUnten}\tag{$\ast\ast$}
%		E:=\set{ε>0:\P(A_ε)>0}\text{ ist höchstens abzählbar},
%	\end{align}
%	denn:
%	\begin{align*}
%		E=\Union_{m  \in \N} \underbrace{\set{ε>0: \P(A_ε) \ge \frac1m}}_{=:E_m}
%	\end{align*}
%	Wir möchten eine Obergrenze für die Größe der $E_m$ finden. Dafür beachte
%	\begin{align*}
%		1\ge \P\klammern{\Union_{ε  \in E_m} A_{ε}}
%		\stackeq{\text{pw.\,disj.}}
%		\sum_{ε  \in E_m}\underbrace{\P(A_{ε})}_{\ge\frac{1}{m}} \ge \abs{E_m}·\frac{1}{m}  \Rightarrow  \abs{E_m} \le m.
%	\end{align*}
%	% Fergers Variante:
%	% Es gilt $\abs{E_m}\le m$, weil: Angenommen es existieren $0<ε_1<…<ε_{m+1}$ mit
%	% % CHECKED: '|'
%	% \begin{align*}
%	% 	&\P(A_{ε_i})\ge\frac{1}{m}\qquad \forall 1\le i\le m+1\\
%	% 	& \Rightarrow 
%	% 	1\ge \P\left(\Union_{i=1}^{m+1} A_{ε_i}\right)
%	% 	\stackeq{\text{pw. disj.}}
%	% 	\sum_{i=1}^{m+1}\underbrace{\P\klammern[\big]{A_{ε_i}}}_{\ge\frac{1}{m}}\ge(m+1)·\frac{1}{m}>1
%	% \end{align*}
%	% Das ist ein Widerspruch.
%	Damit ist $E$ abzählbare Vereinigung endlicher Mengen, also höchstens abzählbar unendlich.
%	Damit liegt das Komplement
%	\begin{align*}
%		E^C=\set[\big]{ε>0: \P(A_ε)=0}
%	\end{align*}
%	dicht in $[0,∞)$.
%	(Dies kann man durch Widerspruch zeigen: in $E$ kann kein Intervall positiver Länge enthalten sein, denn diese enthalten überabzählbar viele Elemente)
%	Daraus folgt insbesondere:
%	\begin{align*}
%		∃(ε_k)_{k \in\N}⊆ℝ\mitε_k\diffskipownarrow0: \forall k \in\N:
%		F_k:=\set[\big]{ x \in\S:d(x,F)\leε_k}\text{ ist $\P$-randlos}
%	\end{align*}
%	Beachte $A_{ε_k}=\partial F_k$. Wähle also $ε_k  \in E^C$ für alle $k  \in \N$. Da $F⊆ F_k~ \forall k \in\N$, gilt:
%	\begin{align*}
%		\limsup_{n \to \infty} \P_n(F)
%		\le\limsup_{n \to \infty} \underbrace{\P_n(F_k)}_{\text{konv.}}
%		\overset{\ref{it:4.2BorelSets}}&{=}
%		\P(F_k)\qquad \forall k \in\N\\
%		\overset{k \to \infty}{ \Rightarrow }
%		\limsup_{n ⟶ ∞} \P_n(F)
%		\le\lim_{k \to \infty} \P(F_k)
%		&=\P(F)
%	\end{align*}
%	Die letzte Gleichheit gilt, weil $\P$ $σ$-stetig von oben ist und $F_k\diffskipownarrow F$.
%	$F_k\diffskipownarrow F$, denn $F_1⊇ F_2⊇…$, da $ε_k$ monoton fallende Folge ist und
%	\begin{align*}
%		\bigcap_{k \in\N}F_k=F,
%	\end{align*}
%	denn:
%	\begin{align*}
%		x \in\bigcap_{k \in\N}F_k
%		&⇔
%		x \in F_k &  \forall k \in\N\\
%		&⇔
%		d(x,F)\leε_k &  \forall k \in\N\\
%		& \Rightarrow 
%		d(x,F)=0\\
%		\overset{\ref{lemma2.3}~\ref{it:distCharakterisierung}}&{⇔}
%		x \in \overline{F}\stackeq{F \in\F}F
%		&& \qedhere
%	\end{align*}
\end{proof}

Falls $\mathcal{U} = \menge{B \in \B(S) : P(\rand{B}) = 0}$, so besagt \cref{satz: 4.2} in \cref{it:4.2BorelSets}, dass $P_n(B) \to P(B)$ für alle $B \in \mathcal{U}$, also punktweise Konvergenz auf der Klasse $\mathcal{U}$, die schwache Konvergenz $P_n \weakto P$ nach sich zieht.
Es gibt weitere solcher Klassen.
% Der letzte Absatz ersetzt folgenden alten Satz:
% Mitunter folgt schwache Konvergenz aus $\P_n(A)\ntoinf  \P(A)$ für eine spezielle Klasse von Mengen $A$.

\begin{theorem} \label{theorem: 4.3}
	Sei $\mathcal{U} \subseteq \B(\S)$ mit
	\begin{enumerate}[label=(\roman*)]
		\item \label{it: 4.3Schnittstabil} $\mathcal{U}$ ist endlich schnittstabil, d.h. mit $A,B \in \mathcal{U}$ ist auch $A \cap B \in \mathcal{U}$
        \item \label{it: 4.3Basis} Jedes offene $G \subseteq \S$ ist abzählbare Vereinigung von Mengen aus $\mathcal{U}$.
	\end{enumerate}
	Dann gilt: $P_n(A) \ntoinf P(A) \enskip \forall A \in \mathcal{U} \follows P_n \weakto P$.
\end{theorem}

\begin{proof}
	Seien $A_1,\dots,A_m \in \mathcal{U}$. Dann gilt:
	\begin{align*}
		P_n \brackets{\bigcup_{i=1}^m A_i}
		&= \sum_{k=1}^m (-1)^{k-1} * \sum_{1 \le i_1 < \dots < i_k \le m} \underbrace{P_n(\overbrace{A_{i_1} \cap \dots \cap A_{i_k}}^{\in\mathcal{U}})}_{\ntoinf  P(A_{i_1}\cap\dots\cap A_{i_k})}\\
		\overset{n\to\infty}&{\longrightarrow}
		\sum_{k=1}^m\sum_{1\le i_1<\dots<i_k\le m} P\brackets{A_{i_1}\cap\dots\cap A_{i_k}}\\
		\overset{\text{allg.\ Add.}}&=
		P\brackets{\bigcup_{i=1}^m A_i}
	\end{align*}
	Sei $G\in\G$, also offen. Wegen Voraussetzung \ref{it: 4.3Basis} ist $G$ darstellbar als $G=\bigcup_{i\in\N} A_i$ mit $(A_i){i \in \N} \subseteq \mathcal{U}$. Es gilt $G_m \defeq \bigcup_{i=1}^m A_i \uparrow G$ für $m\to\infty$. Mit der $\sigma$-Stetigkeit von $P$ gilt
	\begin{equation*}
		\begin{aligned}
			&\forall \epsilon>0 \enskip \exists m_0\in\N \enskip \forall m > m_0: \enskip P(G)-P(G_{m_0})\le\epsilon \\
			\follows &\forall\epsilon > 0 \enskip \exists m_0 \in \N: \enskip
			P(G) - \epsilon \le P(G_{m_0})
			\overset{\text{s.\,o.}}{=}
			\lim_{n \to \infty} P_n(\underbrace{G_{m_0}}_{\subseteq G}) \le \liminf_{n \to \infty} P_n(G) \satzende
		\end{aligned}
	\end{equation*}
	Mit $\epsilon \to 0$ folgt daraus $\liminf_{n \to \infty} P_n(G)\ge P(G)$.
	Da $G \in \G$ beliebig war, folgt nun die Behauptung aus \cref{satz: 4.2} \ref{it:4.2OpenSets}.
\end{proof}

\begin{korollar} \label{korollar4.4}
	Sei $\mathcal{U} \subseteq \B(S)$ endlich durchschnittsstabil mit
	\begin{equation} \label{eqKorollar4.4_i}
		\forall x\in S \enskip \forall \epsilon > 0 \enskip \exists A \in \mathcal{U} : \enskip x \in \inner A \subseteq A \subseteq B(x,\epsilon)
		\satzende
	\end{equation}
	Ist $(\S,d)$ separabel, so gilt
	\begin{align*}
		\brackets{\forall A \in \mathcal{U}: P_n(A) \ntoinf P(A)}
		\follows P_n \weakto P
		\satzende
	\end{align*}
\end{korollar}

\begin{proof}
	Gemäß \cref{satz: 2.9} hat $\G$ eine abzählbare Basis. Nach dem \textit{Satz von Lindelőf} (vergleiche \cite{schubert423277teubner}):
	% wirklich ein Lindelöf mit witzigen ő-Strichen/Punkten
	\begin{center}
		Für jede offene Überdeckung einer beliebigen Teilmenge von $\S$ existiert eine \emph{abzählbare} Teilüberdeckung.
	\end{center}
	Sei nun $G\in\G$ beliebig.
	Für alle $x\in G$ existiert ein $\epsilon_x>0$ mit $B(x,\epsilon_x)\subseteq G$.
	Gemäß \eqref{eqKorollar4.4_i} findet man ein $A_x\in\mathcal{U}$ mit $x\in \inner{A_x}\subseteq A_x\subseteq B(x,\epsilon_x)\subseteq G$.
	Also folgt
	\begin{align*}
		G=\bigcup_{x\in G}\menge{ x}\subseteq\bigcup_{x\in G}\inner{A_x}\subseteq G
	\end{align*}
	Somit ist $\menge{ \inner{A_x}:x\in G}$ eine offene Teilüberdeckung von $G$.
	Aus dem Satz von Lindelőf existieren nun $A_{x_i}\in\mathcal{U}$ ($i\in\N$) mit
	\begin{align*}
		G\subseteq\bigcup_{i\in\N} \inner{A_{x_i}}\subseteq
		\bigcup_{i\in\N}A_{x_i}\subseteq
		G
		\follows
		G=\bigcup_{i\in\N}A_{x_i}
	\end{align*}
	Also erfüllt $\mathcal{U}$ die Voraussetzung \ref{it: 4.3Schnittstabil} und \ref{it: 4.3Basis} in \cref{theorem: 4.3} und es folgt die Behauptung.
\end{proof}

Als Anwendung/ Beschreibung der schwachen Konvergenz im $\S=\R$.
Wir erinnern uns an die \begriff{schwache Konvergenz von Verteilungsfunktionen} $(F_n)_{n\in\N}$ gegen $F$, in Zeichen
\begin{equation*}
	F_n \rightharpoonup F
	\defequiv
	F_n(x) \ntoinf  F(x) \qquad\forall x \in C_F \defeq \menge{ x \in \R: F \text{ ist stetig in }x}
\end{equation*}

\begin{korollar}\label{korollar4.5}
	Seien $P$ und $(P_n)_{n\in\N}$ Wahrscheinlichkeitsmaße auf $(\R, \B(\R))$ mit zugehörigen Verteilungsfunktionen $F$ und $(F_n)_{n\in\N}$ (d.\,h.\ $F(x) = P((-\infty, x])$) für alle $x \in \R$, $n \in \N$).
	Dann gilt:
	\begin{equation*}
		P_n \weakto P
		\equivalent
		F_n \rightharpoonup F
	\end{equation*}
\end{korollar}

\begin{proof}
	\begin{description}
		\hinrichtung
		Sei $x\in\R$ und $B\defeq(-\infty,x]$. Dann ist $\partial B = \menge{x}$ und somit
		\begin{equation}\label{eq:4.5randlos}
			P(\partial B) = 0
			\equivalent P(\menge{ x})=F(x)-\underbrace{F(x-0)}_{\text{Grenzwert}}=0
			\qquad\forall x\in C_F \satzende
		\end{equation}
		Damit folgt für $x\in C_F$
		\begin{equation*}
			F_n(x) \overset{\text{Def}}{=}
			P_n\brackets{\underbrace{(-\infty,x]}_{=B}}
			= P_n(B) \ntoinf  P(B) 
			\overset{\text{Def}}{=} F(x) \qquad\forall x\in C_F
		\end{equation*}
		gemäß \cref{satz: 4.2} \ref{it:4.2BorelSets}, da $B$ $P$-randlos ist für alle $x \in C_F$ gemäß \eqref{eq:4.5randlos}.
		%
		\rueckrichtung
		Das Mengensystem $\mathcal{U} \defeq \menge{ (a,b] : a, b \in C_F \mit a \le b}$ ist endlich durchschnittsstabil. Ferner ist die Menge
		\begin{equation*}
			D_F \defeq C_F^\complement = \menge{ x\in\R: F\text{ \textit{nicht} stetig in }x}
			= \menge{ x\in\R:P(\menge{ x})>0}
		\end{equation*}
		höchstens abzählbar (vgl.\ \eqref{eqProof1.4.2DoppelSternchen}
		im Beweis von \cref{satz: 4.2}). Damit ist $D_F^\complement = C_F$ dicht in $\R$.
		Also folgt
		\begin{align*}
			\forall x\in\R:\forall\epsilon>0:\exists A&=(a,b]\in \mathcal{U}:\\
			x\in (a,b)&=\inner A\subseteq A=(a,b]\subseteq B(x,\epsilon)=(x-\epsilon,x+\epsilon)
		\end{align*}
		denn mit der Dichtheit von $C_F$ gilt:
		In $(x-\epsilon,x)$ muss ein $a\in C_F$ existieren und analog findet man ein $b\in(x,x+\epsilon)$.
		Somit erfüllt $\mathcal{U}$ die Voraussetzungen von \cref{korollar4.4}.
		Klar: $\S=\R$ ist separabel, da $\Q$ abzählbar und dicht in $\R$ ist.
		Schließlich gilt:
		\begin{equation*}
			P_n\brackets{(a,b]} =
			F_n(b)-F_n(a) \overset{a,b \in C_F}{\longrightarrow} F(b)-F(a)
			=P\brackets{(a,b]} \qquad\forall a,b\in C_F
		\end{equation*}
		und mit \cref{korollar4.4} folgt daraus $P_n \weakto P$.
	\end{description}
\end{proof}

\begin{bemerkung} \label{bemerkung4.6} %4.6
	\begin{enumerate}[label=(\arabic*)]
		\item Seien $X$ und $(X_n)_{n\in\N}$ reelle Zufallsvariablen über $(\Omega,\A,\P)$ mit Verteilungsfunktionen $F_n$ bzw.\ $F$. Dann gilt:
			\begin{equation}
				\label{eq:Bemerkung4.6}
				\begin{aligned}
					X_n\distto  X
					\overset{\text{Def}}&{\equivalent}
					P_n = \P\circ X_n^{-1}
					\weakto \P\circ X^{-1} = P \\
					\overset{\text{\cref{korollar4.5}}}&{\equivalent}
					\P(X_n\le x)
					\ntoinf
					\P(X\le x) \qquad \forall x \in C_F \\
					&\equivalent F_n \rightharpoonup F
				\end{aligned}
			\end{equation}
			% für alle $x$, die Stetigkeitsstellen der Verteilungsfunktion von $X$ sind.
			Damit ist das klassische Konzept der Verteilungskonvergenz reeller
			Zufallsvariablen nur ein Spezialfall für $\S = \R$.
		\item \label{it:4.6multiDim} Es gibt Verallgemeinerung von \cref{korollar4.5} bzw \eqref{eq:Bemerkung4.6} auf $\S=\R^k$:
			Seien
			\begin{align*}
				X=\brackets{X^{(1)},\dots, X^{(k)}},
				X_n=\brackets{X^{(1)}_n,\dots,X^{(k)}_n},\qquad(\Omega,\A)\to\brackets{\R^k,\B(\R^k)}
			\end{align*}
			Zufallsvariablen in $\R^k$. Dann gilt:
			\begin{align*}
				X_n\distto  X\text{ in }\R^k
				\equivalent
				\P(X_n\le x)\ntoinf \P(X\le x)\qquad\forall x=\brackets{x_1,\dots,x_k}\in\R^k
			\end{align*}
		wobei $x_i$ Stetigkeitsstelle der Verteilungsfunktion von $X^{(i)}$ ist für alle $i\in\menge{1,\dots,k}$.
		Beweis ist analog zu \cref{korollar4.5}, vergleiche \cite[Satz 5.58]{witting1995mathematische}.
	\end{enumerate}
\end{bemerkung}

Der schwache Limes einer Folge $(P_n)_{n\in\N}$ ist eindeutig.

\begin{lemma}\label{lemma4.6Einhalb}
	$\quad P_n \weakto  P \quad \und \quad P_n \weakto  Q \qquad \follows \qquad P=Q$
\end{lemma}

\begin{proof}
	Gemäß Definition gilt:
	\begin{equation*}
		\int f \diffskip P_n \to \int f \diffskip P 
		\quad \und \quad 
		\int f \diffskip P_n \to \int f \diffskip Q \qquad \forall f\in C^b(\S)
	\end{equation*}
	Der Grenzwert von reellen Zahlenfolgen eindeutig ist, folgt $\int f \diffskip P = \int f \diffskip Q$ für alle $f \in C^b(\S)$ und mit \cref{satz: 3.17} folgt daraus schließlich $P = Q$.
\end{proof}

In \cref{bemerkung4.6} wurde \cref{korollar4.5} über \cref{definition: 4.1} übertragen in eine Version für Verteilungskonvergenz von reellen Zufallsvariablen.
Analog lassen sich auch \cref{satz: 4.2}, \cref{theorem: 4.3} und \cref{korollar4.4} sowie \cref{lemma4.6Einhalb} auf Zufallsvariablen übertragen.

\begin{satz}[Portmanteau-Theorem] \label{satz4.7}
	Folgende Aussagen sind äquivalent:
	\begin{enumerate}[label=(\arabic*)]
		\item $\begin{aligned}
			X_n \distto  X\text{ in }(\S,d)
		\end{aligned}$
		\item $\begin{aligned}
			\EW[f(X_n)] \ntoinf \EW[f(X)] \qquad\forall f \in C^b(\S)
		\end{aligned}$ gleichmäßig stetig
		\item $\begin{aligned}
			\limsup_{n\to\infty} \P(X_n\in F)\le\P(X\in F)\qquad\forall F \in \F
		\end{aligned}$
		\item $\begin{aligned}
			\liminf_{n\to\infty}\P(X_n\in G)\ge\P(X\in G)\qquad\forall G\in\G
		\end{aligned}$
		\item $\begin{aligned}
			\P(X_n\in B)\ntoinf \P(X\in B)\qquad\forall B\in\B(\S)\mit\P(X\in\partial B)=0
		\end{aligned}$
	\end{enumerate}
\end{satz}

\begin{proof}
	Wende \cref{satz: 4.2} an auf $P_n \defeq \P \circ X^{-1}_n$ und $P \defeq \P \circ X^{-1}$.
	Beachte z.\,B.
	\begin{equation*}
		P_n(F) = \P\circ X^{-1}_n(F)
		\overset{\text{Def}}{=}
		\P\brackets{X_n^{-1}(F)}
		=\P\brackets{(\menge{\Omega\in\Omega:X_n(\Omega)\in F}}
		= \P(X_n \in F)
	\end{equation*}
	und
	\begin{equation*}
		\int f \diffskip P_n
		=
		\int_{\S} f\diffskip(\P\circ X_n^{-1})
		\overset{\eqref{eqTrafo}}{=}
		\int_\Omega f(X_n) \diffskip\P
		= \E{f(X_n)}  
		\satzende
	\end{equation*}
\end{proof}

Ferner erhält man unter den jeweiligen Voraussetzungen in \cref{theorem: 4.3} bzw.\ \cref{korollar4.4}:
\begin{equation*}
	\P(X_n\in A) \ntoinf \P(X\in A)\qquad\forall A\in \mathcal{U}
	\qquad \implies \qquad
	X_n \distto  X\text{ in }(\S,d)
\end{equation*}
und aus \cref{lemma4.6Einhalb}:
\begin{equation*}
	X_n\distto  X, \quad X_n\distto  X' 
	\quad \implies \quad
	X \disteq X' \satzende
\end{equation*}

	% This work is licensed under the Creative Commons
% Attribution-NonCommercial-ShareAlike 4.0 International License. To view a copy
% of this license, visit http://creativecommons.org/licenses/by-nc-sa/4.0/ or
% send a letter to Creative Commons, PO Box 1866, Mountain View, CA 94042, USA.

\subsection{Das Continuous Mapping Theorem (CMT)}
Sei $h \colon (\S,d) \to (\S',d')$ messbar.
Wir wollen nun Bedingung an $h$ finden, sodass die beiden folgenden Implikationen gelten:
\begin{enumerate}[label=(\arabic*)]
	\item \label{it: 4.1conditionhweakto} 
	$\begin{aligned}
		P_n \weakto P \follows P_n \circ h^{-1} \weakto P \circ h^{-1}
	\end{aligned}$
	\item \label{it: 4.1conditionhDistrto} 
	$\begin{aligned}
		X_n \distto X \text{ in } (\S,d) \follows h(X_n) \distto  h(X) \text{ in }(\S',d')
	\end{aligned}$
\end{enumerate}
Man beachte dabei, dass
\begin{align*}
	P \circ \brackets{h(X_n)}^{-1}
	&= P \circ \brackets{h \circ X_n}^{-1}
	= \brackets{P \circ X_n^{-1}} \circ h^{-1} \\
	P \circ \brackets{h(X)}^{-1}
	&= \brackets{P \circ X^{-1}} \circ h^{-1}
\end{align*}
gilt. Also folgt \ref{it: 4.1conditionhDistrto} aus \ref{it: 4.1conditionhweakto}. 
Ist $h$ stetig auf $\S$, dann gilt \ref{it: 4.1conditionhweakto} schon. Um das zu zeigen, sei $f\in C^b(\S')$ beliebig --- es gilt:
\begin{align*}
	\int_{\S'} f \diffskip \brackets{P_n \circ h^{-1}}
	\overset{\eqref{eqTrafo}}{=}
	\int_{\S} \underbrace{f \circ h}_{\in C^b(\S)} \diffskip{P_n}
	\overset{\text{Def.}}&{\longrightarrow} \int f \circ h \diffskip P
	\overset{\eqref{eqTrafo}}{=}
	\int f \diffskip{(P \circ h^{-1})} \\
	\overset{\text{Def.}}{\follows}
	P_n \circ h^{-1} &\weakto P \circ h^{-1} \satzende
\end{align*}
Auf Stetigkeit von $h$ kann im Allgemeinen nicht verzichtet werden, denn es gilt das nachfolgende Gegenbeispiel.
%
% Nebenbeibemerkung: Prof. Ferger ist bei research gate angemeldet, das Facebook der Wissenschaftler: dort bekommt er etwas Rückmeldung zu den Papern. Bei üblichen Veröffentlichungen gibt es das nicht. Schade.
%
\begin{beispiel}\label{beispiel: 4.8}
	% Beweis unvollständig, als Übung belassen, hier vollständig
	Sei $\S = \S' = [0,1]$ und $h \colon [0,1] \to [0,1]$ mit
	\begin{align*}
		h(x)\defeq
		\begin{cases}
			1, &\falls x \in \menge{0} \cup \menge{\frac{1}{2n} : n \in \N} \\
			0, & \text{sonst} \satzende
		\end{cases}
	\end{align*}
	Sei $P_n \defeq \delta_{\frac{1}{n}}$ ($n\in\N$) und $P \defeq \delta_{0}$, wobei $\delta_x$ das Dirac-Maß bezeichnet.
	Dann gilt $P_n \weakto P$, denn
	\begin{equation*}
		\int f \diffskip{P_n} = f \brackets{\frac{1}{n}} \ntoinf f(0) = \int f \diffskip P \qquad \forall f \in C^b \brackets{[0,1]} \satzende
	\end{equation*}
	Aber wegen $\delta_x \circ h^{-1} = \delta_{h(x)}$ gilt für
	\begin{align*}
		Q_n \defeq P_n \circ h^{-1} 
		= \delta_{\frac{1}{n}} \circ h^{-1}
		=\delta_{h\brackets{\frac{1}{n}}} 
		\komma \qquad 
		Q \defeq P \circ h^{-1} = \delta_1
	\end{align*}
	das Folgende:
	\begin{equation*}
		\int f \diffskip{Q_n} = f \brackets{h \brackets{\frac{1}{n}}}
		=
		\begin{cases}
			f(1), & \falls n \text{ gerade }\\
			0, & \falls n \text{ ungerade } \satzende
		\end{cases}
	\end{equation*}
	Sei $f \in C^b([0,1])$ mit $f(0) \neq f(1)$ (z.\,B.\ $f = \id$).
	Folglich ist die Folge $\brackets{\int f \diffskip{Q_n}}_{n \in \N}$ divergent.
	Somit ist $\int f \diffskip{Q_n} \to \int f \diffskip{Q}$, was $Q_n \weakto Q$ zeigt.
\end{beispiel}

Die Forderung der Stetigkeit lässt sich aber abschwächen so, dass \cref{it: 4.1conditionhweakto} noch gilt.
Dazu definieren wir die \begriff{Menge der Unstetigkeitsstellen}
\begin{equation*}
	D_h \defeq \menge{ x \in \S : h \text{ nicht stetig in } x } \satzende
\end{equation*}

\begin{lemma} \label{lemma: 4.9}
	Für alle $h \colon \S \to \S'$ gilt $D_h \in \B(\S)$. $h$ muss insbesondere nicht einmal messbar sein.
\end{lemma}

\begin{proof}[nach {\cite[Appendix]{billingsley2013convergence}}]
	Mit der Dreiecksungleichung überlegt man sich leicht:
	\begin{align*}
		h \text{ stetig in } x
		\equivalent
		&\forall 0 < \epsilon \in \Q : \exists 0 < \delta \in \Q : \forall y,z \in \S : \\
		&d(x,y) < \delta \und d(x,z) < \delta
		\follows d'\brackets{h(y),h(z)} < \epsilon
		\satzende
	\end{align*}
	Damit folgt
	\begin{equation} \label{eqProof4.9Plus}
		D_h = \bigcup_{ 0 < \epsilon \in \Q} \bigcap_{0 < \delta \in \Q} \underbrace{\menge{x \in \S :
		\exists y,z \in \S \mit d(x,y) < \delta \und d(x,y) < \delta \und d'\brackets{h(y),h(z)} \ge \epsilon}}_{\defqe A_{\epsilon,\delta}} \satzende
	\end{equation}
	Zeige nun
	\begin{equation}\label{eqProof4.9Stern}
		A_{\epsilon,\delta} \in \G(\S) \qquad \forall \epsilon,\delta > 0 \satzende
	\end{equation}
	Dazu sei $x_0 \in A_{\epsilon,\delta}$.
	Dann existieren $y,z \in \S$ mit $d(x_0,y) < \delta$ und $d(z_0,z) < \delta$, aber $d'\brackets{h(y),h(z)} \ge \epsilon$.
	Wähle
	\begin{equation*}
		\delta_0 \defeq \min\menge{\delta -  d(x_0,y), \delta - d(x_0,z)} > 0 \satzende
	\end{equation*}
	Ferner gilt
	\begin{equation}\label{eqProof4.9i}
		B(x_0,\delta_0) \subseteq A_{\epsilon,\delta},
	\end{equation}
	denn: Sei $x \in B(x_0,\delta_0)$.
	Dann gilt $d(x,y) \le d(x,x_0) + d(x_0,y) < \delta_0 + d(x_0,y) < \delta$
	und analog $d(x,z) < \delta$
	Also folgt $x \in A_{\epsilon,\delta}$, denn $d'\brackets{(h(y),h(z)} \ge \epsilon$.
	Wegen \eqref{eqProof4.9i} ist $x_0$ innerer Punkt von $A_{\epsilon,\delta}$.
	Also folgt \eqref{eqProof4.9Stern} und mit \eqref{eqProof4.9Plus} dann die Behauptung.
\end{proof}

\begin{satz}[Continuous Mapping Theorem -- CMT] \label{satz: 4.10ContinuousMappingTheorem}
	Sei $h \colon (\S,d) \to (\S',d')$ $\B_d(\S)$-$\B_{d'}(\S')$-messbar.
	% hier ist D_h BDSM (B_d(\S) messbar)
	Dann gilt:
	\begin{enumerate}[label=(\arabic*)]
		\item \label{it:4.10CMTweakto} $\begin{aligned}
			P_n \weakto  P \und P(D_h) = 0
			\follows P_n \circ h^{-1} \weakto P \circ h^{-1}
		\end{aligned}$
	\item \label{it:4.10CMTdistrto} $\begin{aligned}
			X_n \distto X \text{ in }(\S,d) \und \P(X \in D_h) = 0
			\follows h(X_n) \distto h(X)\text{ in } (\S',d')
		\end{aligned}$
	\item \label{it:4.10CMThcontinous} Sei $h$ stetig. Dann gilt: $X_n \distto X \text{ in } (\S,d) \follows h(X_n) \distto h(X) \text{ in } (\S',d')$.
	\end{enumerate}
\end{satz}

\begin{proof}
	\textit{Der Beweis ist eine beliebte Prüfungsfrage laut Prof.\ Ferger.} \\
	Wegen des Portmanteau-Theorems (\cref{satz: 4.2}) zeigen wir die äquivalente Aussage über den Limes superior.
	\begin{enumerate}[label=(zu \arabic*), leftmargin=*]
		\item 
		Sei $F\in\F(\S')$ (d.\,h.\ $F\subseteq\S'$ abgeschlossen) beliebig.
		Dann gilt:
		\begin{align}\label{eqProof1.4.10(i)}
			\limsup_{n \to \infty} P_n \circ h^{-1}(F)
			&= \limsup_{n \to \infty} P_n \brackets{\underbrace{h^{-1}(F)}_{\subseteq \abschluss{h^{-1}(F)}}} \\\nonumber
			&\le \limsup_{n \to \infty} P_n \brackets{\underbrace{\abschluss{h^{-1}(F)}}_{\in\F(\S)}} \\
			\overset{\ref{satz: 4.2}}&{\le}
			P\brackets{\abschluss{h^{-1}(F)}}
		\end{align}
		Wir zeigen nun, dass
		\begin{equation} \label{eqProof1.4.10(ii)}
			\abschluss{h^{-1}(F)} \subseteq h^{-1}(F) \cup D_h
		\end{equation}
		gilt. Sei dazu $x \in \quer{h^{-1}(F)}$.
		\begin{enumerate}[label=\textbf{Fall \arabic*:}, leftmargin=*]
			\item Sei $x \in D_h$. Dann folgt daraus schon, dass $x \in h^{-1}(F) \cup D_h$.
			\item Sein nun $x \notin D_h$, d.h. $h$ ist stetig in $x$. Per Definition des Abschlusses existiert eine Folge $(x_n)_{n \in \N} \subseteq h^{-1}(F) \subseteq \S$ mit $x_n \ntoinf x$. Für jede solche Folge gilt
			\begin{equation*}
				\underbrace{h(x_n)}_{\in F} \ntoinf h(x)
				\follows h(x) \in \abschluss{F} = F 
				\follows x \in h^{-1}(F)
				\follows x \in h^{-1}(F) \cup D_h \satzende
			\end{equation*} 
		\end{enumerate}
		Mit \eqref{eqProof1.4.10(ii)} setzt sich die Ungleichungskette \eqref{eqProof1.4.10(i)} fort und wir erhalten
		\begin{equation*}
			P\brackets{\quer{h^{-1}(F)}}
			\overset{\eqref{eqProof1.4.10(ii)}}{\le}
			P\brackets{h^{-1}(F) \cup D_h}
			\le P\brackets{h^{-1}(F)} + \underbrace{P(D_h)}_{= 0}
			= P \circ h^{-1}(F) \satzende
		\end{equation*}
		Schließlich liefern  \eqref{eqProof1.4.10(i)} und \cref{satz: 4.2}\ref{it:4.2ClosedSets}  das gewünschte Resultat $P_n \circ h^{-1} \weakto  P \circ h^{-1}$.
		%
		\item Dies folgt direkt aus \ref{it:4.10CMTweakto} nach Definition \cref{definition: 4.1}, denn es ist
		\begin{align*}
			P \circ (h(X_n))^{-1} = P \circ (h \circ X_n)^{-1} = (P \circ X_n^{-1}) \circ h^{-1}
			&\weakto P \circ X^{-1} \\
			&\weakto (P  \circ X^{-1}) \circ h^{-1} = P \circ (h(X))^{-1} \\
			\overset{\text{Def}}{\follows} h(X_n) &\distto h(X) \satzende
		\end{align*}
		%
		\item Dies folgt wiederum direkt aus \ref{it:4.10CMTdistrto}, da für eine stetige Funktion $h$ die Menge $D_h$ stets leer ist.
	\end{enumerate}
\end{proof}

\begin{bemerkung}[Choque-Kapazität]  
	\label{note: choque_kapazitat}
	Es gibt eine Verallgemeinerung des Wahrscheinlichkeitsbegriffs: die so genannte \begriff{Choque-Kapazität}. Auf der Menge aller Choque-Kapazitätäten lässt sich eine Topologie $\tau_{\text{weak}}$ einführen (die so genannte \emph{schwache Topologie}) und ein entsprechendes Portmanteau-Theorems beweisen, vgl. \cite{Ferger2018WeakCO} und \enquote{Dietmar Ferger (2019): \textsl{unveröffentlicht}}.
\end{bemerkung}
	
    \vfill \hfill $\blacksquare$
    
    \appendix
    \chapter{Grundlagen}
    % This work is licensed under the Creative Commons
% Attribution-NonCommercial-ShareAlike 4.0 International License. To view a copy
% of this license, visit http://creativecommons.org/licenses/by-nc-sa/4.0/ or
% send a letter to Creative Commons, PO Box 1866, Mountain View, CA 94042, USA.

\begin{itemize}[leftmargin=*]
	\item $f: \R \to \R $ heißt \begriff{Borel-messbar}, wenn $f^{-1}(M) \in \B(\R)$ für alle $M \in \B(\R)$.
	\item Sei $(\Omega,\A, \P)$ ein Wahrscheinlichkeitsraum und $(S,\B)$ ein Messraum. Eine \begriff{Zufallsvariable (ZV)} $X$ ist eine $(\A,\B)$-messbare Abbildung $X \colon \Omega \to S$. Wenn $S =\R $, dann heißt $X$ \begriff{Zufallsgröße (ZG)}.
	\item Die \begriff{Verteilung} einer Zufallsgröße $X$ ist Wahrscheinlichkeitsmaß $\mu$ auf $\B$ definiert durch
	\begin{equation*}
		\mu_X \colon \B \to [0,1], \qquad 
		\mu_X(B) \defeq \P(X \in B) \defeq \P(X^{-1}(B)) = \P(\menge{ \omega \in \Omega : X(\omega) \in B}) \satzende
	\end{equation*}
	\item Die \begriff{Dichte} ist
	\begin{align*}
		p \colon \R^n \to [0,\infty] \mit \mu_X(A) = \int_A p(x) \diffskip{\P(x)} \qquad \forall A \in \B(\R^n)
	\end{align*}
	\item \begriff{Verteilungsfunktion} von $X$ ist
	\begin{equation*}
		F_X \colon \R \to [0,1], \qquad F_X(x)=\P(X<x) = \mu_X((-\infty,x)) = \int_{-\infty}^x p(t) \diffskip t
	\end{equation*}
	mit der Eigenschaft $F'(x) = p(x)$ für alle $x \in \R$ fast überall.
	\item Zwei Zufallsvariablen $X_1 \colon \Omega \to S_1$ und $X_2 \colon \Omega \to S_2$ heißen \begriff{unabhängig}, wenn
	\begin{equation*}
		\P(X_1\in B_1) * \P(X_2\in B_2) = \P(X_1\in B_1, X_2\in B_2) 
		= \P\brackets{\menge{\omega \in \Omega : X_1(\omega)\in B_1 \und X_2(\omega) \in B_2}}
	\end{equation*}
	für alle $B_1 \in \B_1$ und alle $B_2 \in \B_2$.
	\item Sei $X$ eine $\P$-integrierbare oder nichtnegative Zufallsgröße. Dann ist der \begriff{Erwartungswert von $X$} definiert als
	\begin{align*}
		\E[X] \defeq \int_\Omega X(\omega) \diffskip\P(\omega) \overset{\text{\ref{eqTrafo}}}{=}
		\int_\R  x \diffskip{\mu_X(x)}
		= \int_\R  x * p(x) \diffskip x
	\end{align*}
	und hat folgende Eigenschaften ($X,Y$ seien Zufallsgrößen):
	\begin{enumerate}
		\item Linearität: $\EW[aX + bY] = a * \EW[X] + b*\EW[Y]$
		\item $X = c \in \R $ fast sicher konstant $\follows \EW[X] = c$
		\item $a \le X \le b$ fast sicher konstant $\follows a \le\E(X) \le b$
		\item $\abs{\EW[X]} \le \EW[\abs{X}]$
		\item $X \ge 0$ fast sicher und $\E(X)=0 \follows X=0$ fast sicher
		\item $X,Y$ unabhängig $\follows \EW[X*Y] = \EW[X] * \EW[Y]$
	\end{enumerate}
	\item Zwei ZG heißen \begriff{unkorreliert}, falls $\E[XY] = \EW[X] * \EW[Y]$.
	\item Für $X \in L^2(\P)$ ist die \begriff{Varianz}
	\begin{equation*}
		\Var(X) \defeq \E[ (X - \E[X])^2] = \int_\Omega (X - \E[X])^2 \diffskip\P = E[X^2] - \E[X]^2
	\end{equation*}
	mit den Eigenschaften
	\begin{align*}
		\Var(a * X + b) &= a^2 * \Var(X) \komma \\
		\Var(X) &=0 \equivalent X \text{ ist konstant fast sicher} \komma \\
		\Var(X + Y) &= \Var(X) + \Var(Y) +\underbrace{2\E[(X-\E[X])*(Y-\E[Y])]}_{=0 \text{, falls $X,Y$ unkorreliert sind}} \satzende
	\end{align*}
\end{itemize}



\begin{satz}[Korrespondenzsatz]\label{satzKorrespondenzsatz}
	Jede Verteilungsfunktion $F$ ist Verteilungsfunktion eines eindeutigen Wahrscheinlichkeitsmaßes $\P$.
	Dieses Maß $\P$ ist durch
	\begin{equation*}
		\P_F((-\infty,x]) \defeq F(x)
	\end{equation*}
	eindeutig bestimmt.
	Umgekehrt bestimmt jedes Wahrscheinlichkeitsmaß eine eindeutige Verteilungsfunktion über
	\begin{equation*}
		F_{\P}(x) \defeq \P((-\infty,x]) \satzende
	\end{equation*}
	Somit ist die Zuordnung der Verteilungsfunktionen zu den Wahrscheinlichkeitsverteilungen bijektiv.
\end{satz}

\begin{bemerkung}[Notation]
	$\P(\diff x) = \diff \P(x)$.
	Außerdem bedeutet $F(\diff x)$ oftmals auch, dass man bzgl. jenem Maß $Q$ integriert, was durch $F$ eindeutig festgelegt ist.
\end{bemerkung}

\begin{satz}[Transformationssatz] \label{satzTransformationssatz}
	Seien $(\Omega,\A,\P)$ ein Wahrscheinlichkeitsraum und $(S,\mathcal{F})$ ein Messraum.
	Sei $g \colon S \to \R $ eine messbare Funktion und $X \colon \Omega \to S$ eine Zufallsvariable.
	Dann gilt:
	\begin{equation} \label{eqTrafo}
		\int_\Omega g(X(\omega)) ~\P(\diff\omega) = \int_S g(s)~\P_X(\diff s)
	\end{equation}
	Hierbei ist $\P_X = \mu_X =\P \circ X^{-1}$ die Verteilung von $X$.
	Für den Standardfall reeller Zahlen ergibt sich mit $f_X$ als Dichte von $X$:
	\begin{equation} \label{eqTrafoR}
		\int_\Omega g(X(\omega))~\P(\diff\omega)
		= \int_\R  g(x)~\P_X(\d x)
		= \int_\R  g(x) * f_X(x) \dx \satzende
	\end{equation}
\end{satz}

\begin{satz}[majorisierte Konvergenz] 
	\label{satzMajorisierteKonvergenz}
	Sei $(\Omega,\A,\P)$ ein Maßraum und sei $(f_n)_{n \in \N}$ eine Folge von $\P$-mesbaren Funktionen (z.\,B.\ ZV) $f_n \colon \Omega \to \R \cup \menge{\infty}$.
	Die Folge $(f_n)_{n \in \N}$ konvergiere $\P$-fast überall gegen eine $\P$-messbare Funktion $f$ und es gilt $\abs{f_n} \le g$ $\P$-fast überall für alle $n \in \N$ für eine $\P$-integrierbare Funktion $g \colon \Omega \to \R $.
	Dann sind $f_n$ und $f$ $\P$-integrierbar und es gilt
	\begin{align*}
		\int_\Omega f(\omega) \diffskip{\P(\omega)} = \int_\Omega \lim_{n \to \infty} f_n(\omega) \diffskip{\P(\omega)}
		= \lim_{n \to \infty}  \int_\Omega f_n(\omega) \diffskip{\P(\omega)}
	\end{align*}
\end{satz}

\begin{satz}[Monotone Konvergenz] \label{satzMonotoneKonvergenz}
	Sei $(\Omega,\A,\P)$ ein Maßraum. Ist $(f_n)_{n \in \N}$ eine Folge nichtnegativer, messbarer Funktionen $f_n \colon \Omega \to [0,\infty]$,
	die $\P$-fast-überall monoton wachsend gegen eine messbare Funktion $f \colon \Omega \to [0,\infty]$ konvergiert, so gilt:
	\begin{align*}
		\int_\Omega f(\omega) \diffskip{\P(\omega)} = \int_\Omega \lim_{n \to \infty} f_n(\omega) \diffskip{\P(\omega)}
		= \lim_{n \to \infty} \int_\Omega f_n(\omega) \diffskip{\P(\omega)}
	\end{align*}
\end{satz}


    
    \nocite{*} % erstellt Literaturverzeichnus ohne Verweise
   	\bibliographystyle{abbrvdin}
    \bibliography{literatur}
\end{document}
